\documentclass[11pt]{article}

\usepackage[margin=1in]{geometry}
\usepackage{amsmath,amssymb,amsthm,mathtools}
\usepackage{mathrsfs}
\usepackage{booktabs}
\usepackage{multirow}
\usepackage{graphicx}
\usepackage{float}
\usepackage{wrapfig}
\usepackage{caption}
\usepackage{subcaption}
\usepackage{enumitem}
\usepackage{xcolor}
\usepackage{microtype}
\usepackage{url}
\usepackage[round]{natbib}
\usepackage[colorlinks=true,citecolor=blue,linkcolor=blue,urlcolor=blue]{hyperref}
\usepackage[nameinlink,noabbrev]{cleveref}

\graphicspath{{figures/}}

\setlist[itemize]{leftmargin=1.2em,itemsep=2pt,topsep=2pt}
\setlist[enumerate]{leftmargin=1.4em,itemsep=2pt,topsep=2pt}

\newcommand{\R}{\mathbb{R}}
\newcommand{\Enc}{E_\theta}
\newcommand{\Dec}{D_\phi}
\newcommand{\shrink}{\operatorname{shrink}}
\newcommand{\blkdiag}{\operatorname{blkdiag}}
\newcommand{\argmin}{\operatorname*{arg\,min}}
\newcommand{\norm}[1]{\left\lVert #1 \right\rVert}
\DeclareMathOperator{\supp}{supp}
\DeclareMathOperator{\sign}{sign}

\theoremstyle{plain}
\newtheorem{proposition}{Proposition}
\newtheorem{theorem}{Theorem}
\newtheorem{lemma}{Lemma}
\newtheorem{corollary}{Corollary}

\theoremstyle{definition}
\newtheorem{definition}{Definition}
\newcommand{\placeholder}[2]{%
  \begin{figure}[tbp]
  \centering
  \fbox{\begin{minipage}{0.92\linewidth}
  \textbf{#1.} #2
  \end{minipage}}
  \caption{#1.}
  \end{figure}
}

\title{Identifying Basins from Sparse Koopman Supports}
\author{Anonymous Authors}
\date{}

\begin{document}
\maketitle

\begin{abstract}
    Koopman autoencoders seek latent coordinates in which nonlinear dynamics evolve linearly, but multibasin systems are challenging: a single finite-dimensional state-reconstructing latent dynamics can hide the basin structure that determines long-run behavior. We show that sparse Koopman autoencoders expose this structure through their active latent supports. The models are trained without basin labels, basin counts, or trajectory-to-basin assignments; basin annotations are used only after training to evaluate the learned supports. Across 15 procedurally generated multibasin systems, merged sparse-support families recover basin-interior structure, with substantially lower conditional basin uncertainty than a dense-latent control. Support interventions further show that the active coordinates are functionally tied to prediction rather than merely correlated with labels. Periodically decoding and re-encoding the current state refreshes support-family selection after controlled basin transfer, and sparse-latent models improve forecasting over dense-latent controls on both the multibasin suite and 10 chaotic Dysts systems. These results identify latent supports as label-free, interpretable basin variables for Koopman learning.
\end{abstract}

\section{Introduction}
\label{sec:introduction}

Learning representations for nonlinear dynamical systems remains a central problem in
machine learning, scientific computing, and control.
A particularly useful goal is to
embed nonlinear dynamics into coordinates in which the evolution is linear, because linear models can be combined with mature tools for prediction,
estimation, and control, including model predictive control and linear quadratic
regulation~\citep{mayne_constrained_2000,grune_nonlinear_2017,korda_linear_2018, kalman_general_1960, mamakoukas_local_2019}. This goal is also well motivated
theoretically. Classical dynamical-systems results provide local topological
linearizations near hyperbolic fixed points~\citep{grobman1959homeomorphism,
hartman1960lemma}, while Koopman eigenfunction constructions can yield linearizing
coordinates on basins of attraction for stable equilibria and periodic
orbits~\citep{lan_linearization_2013,kvalheim_existence_2021}. Thus, nonlinear systems
may admit linear dynamics with a specific change of coordinates.

The difficulty is that these results are primarily existential.
They establish when useful coordinates can exist, but they do not by themselves provide a finite-data procedure for constructing such coordinates.
Koopman operator theory formalizes the linearization idea by lifting nonlinear
state evolution to a linear operator acting on observables \citep{koopman_hamiltonian_1931, koopman_dynamical_1932, mezic2005spectral, brunton_modern_2022}.
This viewpoint has inspired data-driven approximations of Koopman dynamics, such as DMD and eDMD \citep{schmid_dmd_2010, williams_datadriven_2015, williams_kernel-based_2016},
and more recently to Koopman autoencoders, which learn an encoder from state space to a latent representation, a linear latent transition matrix, and a decoder back to state space \citep{takeishi2017learning,lusch_deep_2018,otto2019linearly, azencot_forecasting_2020, shi_deep_2022, fathi2024course}.
These models provide an appealing route from existence results to practical learning: rather than hand-designing observables, they attempt to learn the embedding directly from trajectory data.

A central challenge for Koopman approximations is that many practically interesting nonlinear systems typically contain multiple equilibria or basins of attraction; for example, a damped mechanical system can settle into different wells, a biological or chemical model can approach different stable equilibria, and a controlled system can move between regions with distinct local dynamics. However, theory shows that these systems cannot admit a single finite-dimensional global linearization with a continuous encoder-decoder pair \citep{lan_linearization_2013, brunton_koopman_2016, pan_lifting_2024,liu_properties_2025}, and empirical evidence supports this in practice \citep{fathi2024course}.
In such multistable systems, the representation problem is not only numerical forecasting: a useful model should expose which basin-scale description is active, without being told basin labels during training.
% Given the theoretical argument, multibasin forecasting with a single global finite-dimensional linearization seems ill-posed; it is more naturally a problem of discovering which local or basin-wide linear description applies at the current state.

This paper proposes inducing sparse latent structure as a practical and interpretable method for Koopman learning of multi-attractor systems.
By inducing the state to activate only a small subset of coordinates in lifted latent space, we can incentivize different regions of state space to occupy different active-coordinate patterns.
In this sparse representation, coefficient values remain continuous coordinates for reconstruction and prediction, while the active indices provide a discrete support object that can be inspected across states.

Sparse coding and LISTA-style encoders provide the right inductive bias for this purpose. They naturally induce piecewise-linear, union-of-subspaces structure~\citep{olshausen_emergence_1996,chen_atomic_1998,donoho_optimally_2003,gregor_learning_2010}.
We hypothesize that sparse-latent Koopman autoencoders will assign basin interiors to distinct latent support families, yielding a learned basin variable without basin supervision.
Then, periodically decoding and re-encoding the state can refresh the latent support over time when trajectories move across the state space \citep{fathi2024course}.


% If Koopman models are to serve as a foundation for interpretable multi-regime modeling, then the latent representation of a Koopman autoencoder should do more than forecast well: ideally, it should organize the state space into meaningful local dynamical regimes.
% We hypothesize that methods that induce sparse latents, such as sparse coding, offer an appealing inductive bias for this purpose. Classical sparse coding and LISTA-style encoders produce piecewise-linear mappings and union-of-subspaces structure in latent space \citep{olshausen_emergence_1996, chen_atomic_1998, donoho_optimally_2003, gregor_learning_2010}.
% If the latent representation of state produced by an encoder is sparse, then each state activates only a small subset of latent coordinates.


\paragraph{Contributions.} This paper contributes an empirical and mechanistic study showing (1) induced latent sparsity produces support families that are basin-aligned without basin supervision, (2) active-coordinate interventions show that these supports are functionally used by the learned predictor, (3) periodic re-encoding updates \citep{fathi2024course} refresh support-family selection after controlled basin transfer, and (4) sparse-latent Koopman autoencoders improve long-horizon forecasting performance on both multibasin and chaotic-flow benchmarks.
The results demonstrate that basin identification can be an \textit{emergent} property of sparse-latent Koopman autoencoders.

\section{Problem formulation}
\label{sec:problem}

The introduction motivates a multibasin version of Koopman learning. Given sampled trajectories of a nonlinear system, the aim is to learn a state-reconstructing latent representation whose evolution is linear, without being told which basin or local dynamical regime generated each observation.
% The benchmarks may be defined by continuous-time dynamics, but the learner is only given stored observations.
We formulate the problem at the cadence of the stored benchmark observations.
For each system, let \(\mathcal X\subseteq\R^{d_x}\) be the observed state space and let
\begin{equation}
    x_{k+1}=F_{\Delta t}(x_k),\qquad x_k\in\mathcal X,
    \label{eq:sampled_map}
\end{equation}
where $x_k\in\R^{d_x}$ is the state at the $k$th stored sample, $F_{\Delta t}$ denotes the transformation to the next discrete observation step based on underlying dynamics, and $\Delta t$ is the system-specific observation interval.
% The model does not observe the continuous flow, the vector field, or the internal integration steps used to generate the cached trajectories.
Forecasting is therefore defined by repeatedly applying the map $F_{\Delta t}$; a horizon of $h$ means $h$ applications of $F_{\Delta t}$; we denote $k$ compositions as $F_{\Delta t}^{\,(k)}$. Because $\Delta t$ may differ across systems, equal values of $h$ may yield different physical time.

\paragraph{Basins of attraction.} Many systems of interest are multistable: different initial conditions can approach different long-run behaviors.
A fixed point is a state $x^\ast$ satisfying $F_{\Delta t}(x^\ast)=x^\ast$.
More generally, a trajectory may approach an attractor $A$, such as a stable equilibrium, limit cycle, countable finite set, or chaotic invariant set.
The basin of attraction of $A$ is the set of initial conditions whose trajectories converge to $A$:
\begin{equation}
  \mathcal{B}(A)
  =
  \left\{
    x_0\in\R^{d_x}:
    \inf_{a\in A}\norm{F_{\Delta t}^{\,(k)}(x_0)-a}_2\!\to 0
    \text{ as } k\to\infty
  \right\}.
  \label{eq:basin}
\end{equation}
When multiple attractors coexist, their basins partition the state space up to basin boundaries. This defines the setup for this paper.
% Within one basin, a useful linear description may exist locally or even basin-wide in appropriate coordinates, but another basin may require an incompatible description. Thus,
The task is to learn a representation that can predict while preserving inspectable information about where the state lies in the multibasin geometry.

\paragraph{Koopman theory.} The stored-step map \(F_{\Delta t}\) induces a Koopman operator \(\mathcal K\) on scalar observables \(g:\mathcal X\to\R\) by composition, \((\mathcal K g)(x)=g(F_{\Delta t}(x))\). Although \(F_{\Delta t}\) may be nonlinear, \(\mathcal K\) is linear on the space of observables. A finite-dimensional Koopman approximation searches for a vector observable \(\Phi:\mathcal X\to\R^{d_z}\) and a matrix \(K\) such that \(\Phi(F_{\Delta t}(x))\approx K\Phi(x)\), while retaining enough information to reconstruct the state. Additional background is provided in \cref{app:background}.


\paragraph{Koopman autoencoders.} We study this setting with Koopman autoencoders (KAEs). The model learns an encoder $\Enc:\R^{d_x}\rightarrow\R^{d_z}$, a decoder $\Dec:\R^{d_z}\rightarrow\R^{d_x}$, and a latent linear transition matrix $K\in\R^{d_z\times d_z}$. For observation $x_k$ at arbitrary timestep $k$, the corresponding open-loop $h$-step forecast is
\begin{equation}
  \hat x_{k+h}=\Dec\!\left(K^h\Enc(x_k)\right).
  \label{eq:kae_forecast}
\end{equation}
The same matrix $K$ is applied at every step of the rollout, so any state-dependent structure required for prediction must be encoded in the learned representation rather than in a supplied regime label.

\paragraph{Training windows.} Training uses windows of adjacent stored observations. For a batch of $B$ windows of rollout length $L$,
$\{x_{b,k},x_{b,k+1},\ldots,x_{b,k+L}\}_{b=1}^{B}$, where $b$ indexes windows and $\ell$ indexes offsets within a window, the model encodes each observed state, rolls out from the initial encoded state, and decodes both reconstructed observations and latent forecasts. The encoded/reconstructed quantities are defined for $\ell=0,\ldots,L$, while rollout/forecast quantities are defined for $\ell=1,\ldots,L$:
\begin{equation}
\begin{aligned}
  z_{b,k+\ell}^{\rm enc} = \Enc(x_{b,k+\ell}),\quad
  z_{b,k+\ell}^{\rm roll}=K^\ell z_{b,k}^{\rm enc},\quad
  x_{b,k+\ell}^{\rm rec} = \Dec(z_{b,k+\ell}^{\rm enc}),\quad
  \hat x_{b,k+\ell}=\Dec(z_{b,k+\ell}^{\rm roll}).
\end{aligned}
\label{eq:window_rollout}
\end{equation}
The training losses described later are built from these reconstructed and rolled-out quantities. Each application of $K$ is interpreted as one stored benchmark step, matching the observation map in \cref{eq:sampled_map}.

\paragraph{Model discovers basins unsupervised.} In the intended training and deployment setting, the sparse-support mechanism is not given basin labels, the number of basins, or trajectory-to-basin assignments. When benchmark basin labels are available, they are reserved for post-hoc evaluation; they are not used to choose support families, train the encoder or transition, or select supports. The block-diagonal and soft-block transition rows are architecture diagnostics: for those ablations only, the fixed number of latent transition groups is set from benchmark generator metadata, as detailed in \cref{app:experimental_details}.

\section{Method}
\label{sec:method}
\subsection{Sparse supports as local regime variables}
\label{sec:method_regime_supports}

 A finite-dimensional Koopman representation of the multibasin systems described in \cref{sec:problem} should do more than minimize rollout error. It should also make the basin-scale organization of the state space inspectable. If the representation of a state \(x\) is a sparse code \(z\), then the support, or set of nonzero coordinates, is a discrete object that can be compared across states. The continuous values of \(z\) carry state information, while the support can act as a model-produced candidate basin variable.

 % We call such a discrete indication a \emph{local regime variable}: a quantity computed from the model state that can group states or route predictors without being supplied as a label.


Given an overcomplete dictionary $D$, the Lasso loss \citep{beck2009fast} finds the fewest nonzero coefficients of $z$ to reconstruct the input $x$. In other words, it finds the sparsest solution to the underdetermined problem.
\begin{equation}
  z^\star(x)
  =
  \argmin_{z\in\R^{d_z}}
  \frac{1}{2}\norm{x-Dz}_2^2+\lambda_{\rm sc}\norm{z}_1,
  \qquad \lambda_{\rm sc}>0.
  \label{eq:sparse_coding_objective}
\end{equation}
The classical solver to this problem (ISTA \citep{beck2009fast}) can be unrolled into a deep network of depth $L$, where each loop (layer) corresponds to a single iteration of ISTA. This is known as learned ISTA or LISTA \citep{gregor_learning_2010}, and the sparse code is given by the encoder \(z=\Enc(x)\) where the \emph{support} is the set of active latent coordinates. With a linear decoder with unit-norm columns, the support selects which decoder columns participate in reconstructing the state, while the nonzero coefficients give coordinates within that selected reconstruction subspace. If both $D$ and $z$ are optimized jointly \citep{olshausen_emergence_1996}, the input space is represented through a learned collection of active dictionary-atom subsets. We therefore posit that sparse Koopman autoencoders can capture multibasin attractor structure without explicit labels: the support set identifies the active basin-scale family, and the coefficient values carry within-family state information.


We therefore evaluate the following chain:
\begin{equation}
  \text{sparse latent code}
  \Longrightarrow
  \text{support family--basin alignment}
  \Longrightarrow
  \text{functional and refreshable support selection}.
  \label{eq:claim_chain}
\end{equation}
The sparse latent code is induced by the architecture and objective. Basin labels, when available in benchmarks, are used only after training to evaluate alignment or define diagnostic interventions.

\subsection{Koopman autoencoder and sparse encoder}
\label{sec:method_model}

The predictor follows the Koopman-autoencoder formulation in \cref{sec:problem}: an encoder \(\Enc\), a decoder \(\Dec\), and a latent transition \(K\) produce the rollout in \cref{eq:kae_forecast}. For the linear-decoder rows used here, the learned parameter tuple is \((\theta,D,K)\): encoder parameters \(\theta\), a decoder dictionary \(D\) with normalized columns, and the latent transition. Thus \(\Dec(z)=Dz\), sparsity is imposed on the learned latent representation \(z_t=\Enc(x_t)\), and we vary whether \(K\) is dense, block diagonal, or softly block-regularized. Sparsity is not an assumption about the physical state \(x_t\).

\paragraph{Sparse encoding.} Our encoders implement the sparse-code bias through shrinkage and a training $\ell_1$ sparsity penalty rather than by solving \cref{eq:sparse_coding_objective} at evaluation time. Our main sparse encoder is a learned iterative shrinkage-thresholding architecture, or LISTA encoder \citep{gregor_learning_2010}. It first computes a dense pre-code \(c_t=f_\theta(x_t)\), then applies learned shrinkage refinements for $q=0,\ldots,Q-1$:
\begin{equation}
  u_t^{(0)}=\shrink(c_t,\tau),\qquad
  u_t^{(q+1)}=\shrink(Su_t^{(q)}+c_t,\tau),\qquad
  z_t=u_t^{(Q)}.
  \label{eq:lista}
\end{equation}
Here $\shrink(a,\tau)=\operatorname{sign}(a)\max(|a|-\tau,0)$
is applied elementwise, \(S\) is learned, and \(\tau\) controls the threshold scale for shrinkage. The thresholding step creates exact zeros, so the encoder explicitly selects active latent coordinates while learning the selection rule from the prediction objective.

\paragraph{Transition families.} We vary the structure of \(K\) separately from the encoder. The dense transition leaves \(K\) unrestricted. The block-diagonal transition partitions latent coordinates into \(M\) fixed groups and constrains
\begin{equation}
  K_{\rm block}=\blkdiag(K_1,\ldots,K_M).
  \label{eq:block_transition}
\end{equation}
We also use a soft-block LISTA ablation that keeps \(K\) dense but penalizes cross-group entries,
\begin{equation}
  \mathcal{L}_{\rm offblock}
  =
  \sum_{i,j:\,g(i)\ne g(j)} |K_{ij}|,
  \label{eq:offblock_penalty}
\end{equation}
where \(g(i)\) is the fixed architectural group containing coordinate \(i\).

\subsection{Training objective}
\label{sec:method_training}

Training uses adjacent observation windows. For a batch of $B$ windows with prediction horizon $L$, the model produces reconstructed future states $x_{b,k+\ell}^{\rm rec}$, latent rollouts $z_{b,k+\ell}^{\rm roll}$, encoded future latents $z_{b,k+\ell}^{\rm enc}$, and decoded rollout predictions $\hat x_{b,k+\ell}$ as defined in \cref{eq:window_rollout}. The rollout is seeded by the initial state, and the following sums are over offsets $\ell=1,\ldots,L$:
\begin{align}
  \bar{\mathcal{L}}_{\rm pred}
  &=\frac{1}{BL\sqrt{d_x}}\sum_{b,\ell}\norm{\hat{x}_{b,k+\ell}-x_{b,k+\ell}}_2,
  &
  \bar{\mathcal{L}}_{\rm rec}
  &=\frac{1}{BL\sqrt{d_x}}\sum_{b,\ell}\norm{x_{b,k+\ell}^{\rm rec}-x_{b,k+\ell}}_2,
  \label{eq:loss_state_terms}
  \\
  \bar{\mathcal{L}}_{\rm lin}
  &=\frac{1}{BL}\sum_{b,\ell}\norm{z_{b,k+\ell}^{\rm roll}-z_{b,k+\ell}^{\rm enc}}_2,
  &
  \bar{\mathcal{L}}_{\rm sp}
  &=\frac{1}{BL}\sum_{b,\ell}\norm{z_{b,k+\ell}^{\rm roll}}_1 .
  \label{eq:loss_latent_terms}
\end{align}
The observation-space terms are normalized by $\sqrt{d_x}$ so that their scale is comparable across systems of different state dimension. The latent terms are left in their native scale because the latent dimension is fixed within each comparison.

The total objective is
\begin{equation}
  \mathcal{L}
  =
  \frac{1}{L}\left(
  \lambda_{\rm pred}\bar{\mathcal{L}}_{\rm pred}
  +\lambda_{\rm rec}\bar{\mathcal{L}}_{\rm rec}
  +\lambda_{\rm lin}\bar{\mathcal{L}}_{\rm lin}
  +\lambda_{\rm sp}\bar{\mathcal{L}}_{\rm sp}
  \right)
  +\mathcal{L}_{\rm struct}.
  \label{eq:training_objective}
\end{equation}
Here $\mathcal{L}_{\rm struct}=0$ for the dense and block-diagonal transition variants.
For the soft-block transition, $\mathcal{L}_{\rm struct}=\lambda_{\rm off}\mathcal{L}_{\rm offblock}$, which probes whether a soft transition grouping is sufficient to induce useful support structure.
The prediction and latent-linearity terms make the representation useful for Koopman rollout, the reconstruction term keeps the code tied to the observed state, and the sparsity term encourages selective activation.

% Compared with the minimal Koopman autoencoder objective in \cref{eq:KAE-standard-objective}, this rollout objective adds decoded multi-step prediction, latent consistency over stored observation windows, explicit sparsity pressure, and any transition-structure penalty. These additions are needed because the paper is not only asking whether one-step latent features can be made approximately linear; it is asking whether the learned support remains useful for finite-horizon prediction and can be inspected as a label-free local-regime object.

\subsection{Support objects}
\label{sec:method_supports}

Given \(z=\Enc(x)\in\mathbb{R}^{d_z}\), a \textit{support rule} first converts the latent code into a Boolean active-coordinate vector \(m(x)\in\{0,1\}^{d_z}\), where \(m_i(x)=1\) means that latent coordinate \(i\) is active for state \(x\). Equivalently, the exact support is the index set \(S(x)=\{i:m_i(x)=1\}\). Exact supports defined by a threshold can be too fine-grained: weak coordinates may cross the threshold, and a single basin may be covered by several nearby or partially overlapping active-coordinate vectors.

We therefore distinguish exact supports from \emph{support families}, which group nearby Boolean support vectors and test whether these groups form basin-scale objects.
We use three support constructions:
\begin{itemize}
  \item \textbf{Absolute-threshold support.} The absolute support vector is the Boolean mask
  \[
    m_{\rm abs}(x)=\bigl(\mathbf{1}\{|z_i|>10^{-3}\}\bigr)_{i=1}^{d_z}\in\{0,1\}^{d_z},
    \qquad
    S_{\rm abs}(x)=\{i:m_{{\rm abs},i}(x)=1\}.
  \]
  Here \(m_{{\rm abs},i}(x)\) is the \(i\)-th entry of the mask, and \(S_{\rm abs}(x)\) is the equivalent active-index set. This strict active-coordinate mask is the input to \(F_{\rm abs}\) and to canonical wrong-support interventions.
  % Its entropy, fragmentation, and active-coordinate count are appendix checks.
  \item \textbf{Absolute-threshold support family.} \(F_{\rm abs}\) groups exact supports whose active-index sets substantially overlap. The overlap score is the Jaccard similarity, or intersection-over-union, between the active coordinates of two Boolean support vectors \citep{jaccard_etude_1901}. On an evaluation collection \(\mathcal X\), the algorithm first counts all distinct Boolean vectors and then visits those distinct vectors from most to least frequent, following a leader-style threshold clustering rule \citep{hartigan_clustering_1975}. Each family \(f\) stores one fixed representative vector \(r_f\): the exact support vector that created that family. For the next distinct vector \(m\), the algorithm computes
  \[
    J(m,r_f)=
    \frac{\sum_i \mathbf{1}\{m_i=1\ \text{and}\ r_{f,i}=1\}}
         {\sum_i \mathbf{1}\{m_i=1\ \text{or}\ r_{f,i}=1\}},
  \]
  with \(J(0,0)=1\) when both vectors are all zero. If the best existing representative has Jaccard similarity at least \(0.5\), every occurrence of \(m\) receives that family label. Otherwise \(m\) starts a new family and becomes the representative \(r_f\) for that family. The representatives are not learned centroids and are not updated after assignment; they are exact support vectors chosen by each deterministic frequency-ordered pass. Thus the procedure is closer to threshold-based leader clustering than to \(k\)-means: the number of families is induced by the threshold and the data order, not fixed in advance. We write \(F_{\rm abs}(x)\) for the family assigned to \(m_{\rm abs}(x)\), and \(|F_{\rm abs}|\) for the number of occupied families.
  % Families are built without basin labels; labels are used only afterward to compute \(H(B\mid F_{\rm abs})\).
  \item \textbf{Fixed-size top-\(8\) support and family.} \(S_{\rm top8}(x)=\operatorname{top}_8(|z|)\) is the set of the eight largest-magnitude coordinates. Unlike \(S_{\rm abs}\), it gives every state the same number of active indices and does not depend on a magnitude threshold. \(F_{\rm top8}(x)\) is obtained by applying the same Jaccard merge to exact \(S_{\rm top8}\) support vectors. Each state still begins with exactly eight active indices; \(F_{\rm top8}\) is a family label for nearby eight-index support vectors, not a new fixed-size support.
\end{itemize}
The main basin-alignment object is \(F_{\rm abs}\), which preserves the thresholded active set and then merges nearby exact masks. The support-refresh diagnostic uses \(F_{\rm top8}\), whose underlying \(S_{\rm top8}\) support gives every state the same number of active indices and is invariant to global rescalings of \(z\) when there are no top-\(8\) ties. A scale-relative rule, $S_{\rm rel}(x)=\{i:\ |z_i|>0.1\max_j |z_j|\}$, is used only in appendix sensitivity checks.


\subsection{Periodic re-encoding for support refresh}
\label{sec:method_refresh}
If a support indicates the current basin-scale representation, a long rollout should be able to update that support as the predicted state changes. We therefore evaluate periodic re-encoding, following \citet{fathi2024course}. Starting from $z_k=\Enc(x_k)$ and a re-encoding period $m$, the model advances for $m$ predicted steps under the learned Koopman transition, decodes the predicted state, and encodes that prediction again:
\begin{equation}
  \tilde z_{k+m}=K^m z_k,\qquad
  \tilde x_{k+m}=\Dec(\tilde z_{k+m}),\qquad
  z_{k+m}=\Enc(\tilde x_{k+m}).
  \label{eq:reencode}
\end{equation}
The same procedure is repeated over the rollout horizon. Each rollout uses only the model's predicted state; it does not use the true future state, a basin label, or an oracle for transitions between basins. For the controlled support-refresh diagnostic in \cref{sec:experiments}, we isolate the representation mechanism: the latent is carried by the learned global \(K\) between scheduled events, and the stale pre-refresh \(F_{\rm top8}\) family is compared with the family obtained by encoding the current controlled transfer state. No support-conditioned local map or support-gated transition is fit or used.


\section{Experimental procedure}
\label{sec:experiments}
The experiments ask whether sparse Koopman autoencoders learn basin-identifying supports without basin supervision, and whether those supports remain useful under prediction. We evaluate (i) forecasting on controlled multibasin systems and external chaotic flows, (ii) support-family alignment with withheld basin labels, (iii) wrong-support interventions that test functional use of active coordinates, and (iv) periodic support refresh after controlled basin transfer.

\begin{figure}[tbp!]
\centering
\includegraphics[width=0.8\linewidth]{neurips_paper_2026/fig_benchmark_support_dysts_composite.pdf}
\caption{Benchmark overview. Top: procedurally generated multibasin systems where gray streamlines are the true vector field, colored dots mark support-family/basin matches, and black dots mark support-family/basin mismatches. Learned LISTA support families are mapped post hoc to each family's dominant evaluation basin only for visualization. Bottom: $dt{\times}30$ Dysts horizon-5000 phase portraits with held-out truth in gray and forecasts from the best seed-$0$ LISTA variant.}
\label{fig:benchmarks}
\end{figure}

\paragraph{Benchmarks and training.}
The main benchmark for our interpretability experiments contains 15 procedurally generated two-dimensional multibasin systems, including multiwell potentials, gated vector fields, depth cascades, and polygonal and transition-rich layouts. Basin labels are available only for evaluating alignment and diagnostic metrics. We also use a subset of 10 chaotic systems from the Dysts repository \citep{gilpin_chaos_2021,gilpin_model_2023} with the step size $dt$ multiplied by a factor of 30 as an additional forecasting test.
Analytic expressions and more details are in \cref{app:fixed17_inventory,app:dysts_inventory_full}.

We compare the six model rows shown in \cref{tab:forecasting_main}: three LISTA sparse encoders (where LISTA-BD denotes block diagonal latent transitions and LISTA-SB denotes soft-block regularization),
two sparse-latent MLP controls applying the sparsity penalty $\bar{\mathcal{L}}_{\rm sp}$ \eqref{eq:loss_latent_terms} with dense \citep{fathi2024course} or block-diagonal (BD) transitions,
and a dense-latent MLP control with no sparsity incentive.
Each model is trained for 15 random seeds on each system. Exact training details are reported in \cref{app:experimental_details}.

\paragraph{Forecasting and statistics.}
\label{sec:experiments_forecasting}
\label{sec:stats_protocol}
Forecasting is evaluated on held-out trajectories at the stored observation cadence and summarized by mean squared error (MSE) at prespecified horizons. Forecasting results report the best value over a fixed set of re-encoding periods (in Appendix \cref{app:experimental_details}) unless otherwise stated, where the rollout uses only predicted states and never the true future trajectory.

Point estimates are interquartile means (IQM) over seeds within each system, then arithmetic means across systems. We use seed IQMs within each system rather than the arithmetic mean because few-run deep-learning evaluations are sensitive to outlier runs, and IQM is a robust aggregate that reduces uncertainty in benchmark comparisons \citep{agarwal2021deep}.
Horizon-trend lines use the same point estimate as the tables, with seed-bootstrap bands that keep the benchmark systems fixed and summarize uncertainty over seeds for an average system.
% Specifically, each bootstrap seed-IQM draw is converted to a log-relative perturbation from that system's observed seed-IQM, the perturbations are averaged across systems, and the resulting percentile interval is exponentiated and anchored to the displayed arithmetic-mean line.
Details on system-by-system performance are reported in \cref{app:fixed17_full,app:dysts_full}; statistical tests are described in \cref{app:statistical_testing}.
\paragraph{Support alignment.} This mechanistic experiment examines the extent to which sparse support families fit after training can behave like basin-interior regime variables.
If a support family indexes the local chart or basin-specific description being used by the model, then on states deep inside a basin, it should leave little uncertainty about the withheld benchmark basin label.
States near separatrices are challenging and may confound the inferences: near these states, support changes, shared supports, and larger uncertainty would be mechanistically appropriate.
We therefore evaluate this diagnostic on held-out states deep in basins because basin interiors are the cleanest population for testing whether support items can identify basins.

For a state \(x\), let \(d_1(x)\) and \(d_2(x)\) be its distances to the nearest and second-nearest benchmark attractor centers. The \textit{basin-depth margin} is \(d_2(x)-d_1(x)\).
The support-basin alignment experiment uses ground-truth benchmark geometry only for evaluation, taking the top quartile of held-out states ranked by basin-depth margin within each benchmark basin.
The main metrics are conditional entropy \(H(B\mid F_{\rm abs})\), which should be small when support families are basin-aligned, and \(|F_{\rm abs}|\), a system family-count diagnostic compared with the number of basins represented.

% Conditional entropy is calculated as
% $$H(B\mid F_{\rm abs})
%   = -\sum_{f\in\mathcal F_{\rm abs}} \frac{n_f}{N}
%   \sum_{b\in\mathcal B} \frac{n_{bf}}{n_f}
%   \log \frac{n_{bf}}{n_f},$$
% where $n_{bf}$ is the number of evaluated states with basin $b$ and support family $f$, $n_f=\sum_b n_{bf}$, and $N=\sum_f n_f$.

Functional use is tested by wrong-support and coordinate-intervention ablations.
For the cross-system wrong-support ablation, we construct a canonical deep-slice \(S_{\rm abs}\) mask per basin. Systems whose deep slice contains only one basin are excluded. Then, a held-out rollout replaces the inferred mask with a canonical mask from another basin and holds it fixed. The reported value is the wrong-support/base error ratio.
For a representative sparse model, we also run a more direct inference-time coordinate intervention that separates which sparse coordinates are active from the continuous values placed on those coordinates. Starting from \(z_0=\Enc(x_0)\), we perturb only \(z_0\), then use the ordinary learned latent transition and decoder for the full autonomous rollout. In the coordinate-dropping intervention, the latent entries \(z_{0,i}\) with \(i\in S_{\rm abs}(x_0)\) are ranked by \(|z_{0,i}|\), and the largest entries are cumulatively zeroed. In the random-support intervention, the active coefficient values are preserved but reassigned to randomly chosen inactive latent coordinates. No model parameters are retrained, and the true future trajectory is used only for scoring.

\paragraph{Periodic re-encoding for support refresh.} This experiment examines whether periodic re-encoding refreshes support-family identity in different regions of state space. It evaluates \(F_{\rm top8}\) selection after controlled transfer into a new basin. The latent state is advanced with the learned global \(K\) between scheduled refresh events; at each event, the stale family implied by the propagated latent is compared with the family obtained by re-encoding the current controlled transfer state. Basin labels are used only to construct and score these transfers, and no support-conditioned rollout is used. The reported quantities are the target-family fraction before and after re-encoding at each period, plus a no-transfer false-target control.
Additional procedural details are collected in \cref{app:experimental_details,app:statistical_testing}.

\section{Results}
\label{sec:results}

The results support three main conclusions. Sparse-latent models improve forecasting over dense-latent controls on both benchmark families; sparse supports expose basin-scale support families \(F_{\rm abs}\); and support families can be refreshed by periodic re-encoding after controlled basin transfer.

\begin{table}[tbp]
\centering
\caption{Forecasting performance on the two benchmark families. Values are MSEs; lower is better. Each cell is an arithmetic mean across systems after summarizing seeds within each system by IQM. Superscripts mark Holm-corrected tests against the dense-latent MLP baseline: multibasin cells use within-system Wilcoxon/Holm tests and Dysts cells use exact sign tests across systems (\(\ast\), \(p<0.05\)). \textbf{Bold} marks the best entry in each column.}
\label{tab:forecasting_main}
\label{tab:dysts_long}
% May 6 status: Sparse MLP-BD cells are provisional. Repaired GenericKM block-diagonal reruns are live again under resume jobs 9488208/9488210--9488212 and 9488209/9488213.
{\setlength{\tabcolsep}{2pt}
\resizebox{\textwidth}{!}{\begin{tabular}[t]{@{}l rrr rrr@{}}
\toprule
& \multicolumn{3}{c}{Multibasin, 15 systems} & \multicolumn{3}{c}{Dysts \(dt{\times}30\), 10 systems} \\
\cmidrule(lr){2-4}\cmidrule(l){5-7}
Model & H100 & H500 & H1000 & H100 & H2000 & H4000 \\
\midrule
LISTA & ${0.0413}^{\ast}$ & ${0.154}^{\ast}$ & ${0.176}^{\ast}$ & ${2.65{\times}10^{-4}}^{\ast}$ & $\mathbf{0.0999}$ & ${0.465}^{\ast}$ \\
LISTA-BD & ${0.0414}^{\ast}$ & ${0.126}^{\ast}$ & ${0.147}^{\ast}$ & $3.69{\times}10^{-4}$ & ${0.122}^{\ast}$ & ${0.492}^{\ast}$ \\
LISTA-SB & ${\mathbf{0.0387}}^{\ast}$ & ${0.125}^{\ast}$ & ${0.150}^{\ast}$ & $5.62{\times}10^{-4}$ & $0.182$ & $0.744$ \\
Sparse MLP, BD & ${0.0477}^{\ast}$ & ${0.150}^{\ast}$ & ${0.178}^{\ast}$ & $5.06{\times}10^{-4}$ & $0.130$ & ${\mathbf{0.441}}^{\ast}$ \\
Sparse MLP \citep{fathi2024course} & ${0.0400}^{\ast}$ & ${\mathbf{0.0955}}^{\ast}$ & ${\mathbf{0.109}}^{\ast}$ & ${\mathbf{2.38{\times}10^{-4}}}^{\ast}$ & ${0.122}^{\ast}$ & $0.506$ \\
\midrule
Dense MLP \citep{lusch_deep_2018} \emph{[baseline]} & $0.832$ & $2.68$ & $2.91$ & $0.00125$ & $0.228$ & $0.767$ \\
\midrule
DMD & $2.58$ & $2.14$ & $1.70$ & $2.85$ & $4.19$ & $3.83$ \\
Polynomial EDMD & $0.864$ & $0.968$ & $0.887$ & $2.32$ & $3.93$ & $3.61$ \\
RBF-dictionary EDMD & $0.660$ & $0.808$ & $0.751$ & $0.890$ & $2.84$ & $2.91$ \\
$k$-means local linear & $1.40$ & $3.04$ & $138$ & $2.29$ & $1.5{\times}10^{150}$ & $3.2{\times}10^{140}$ \\
GMM local linear & $1.65$ & $3.57$ & $3.93$ & $10.6$ & $7.8{\times}10^{18}$ & $4.5{\times}10^{44}$ \\
Soft-gated GMM local linear & $1.66$ & $3.71$ & $4.07$ & $5.77$ & $3.0{\times}10^{13}$ & $4.1{\times}10^{34}$ \\
\bottomrule
\end{tabular}
}}
\end{table}

\paragraph{Sparse-latent models improve forecasting (\cref{tab:forecasting_main}).}
On the 15 multibasin systems, all sparse-latent rows reduce MSE relative to the dense-latent MLP at every reported horizon, with reproducible within-system improvements in all 15 systems. The advantage is large: at \(H1000\), the best sparse-latent row has MSE \(0.107\) versus \(2.93\) for Dense MLP. On Dysts, sparse-latent rows also improve over the dense-latent baseline at the displayed representative horizons, with LISTA giving the best longer-horizon values at \(H2000\) and \(H5000\).

\begin{table}[tbp]
\centering
\caption{Support diagnostics and periodic support refresh. \textbf{Left:} basin-support diagnostics on the per-basin deep slice of the 15 multibasin systems. \(H(B\mid F_{\rm abs})\) is conditional basin uncertainty given the learned support family; \(\overline{|F_{\rm abs}|}\) is the mean occupied support-family count; wrong-support ratios compare rollouts using an incorrect canonical \(S_{\rm abs}\) mask with the base rollout. \textbf{Right:} controlled-transfer \(F_{\rm top8}\) support-family refresh using only the learned global \(K\). Each period cell shows the target-family fraction before re-encoding (stale propagated support family) \(\to\) after re-encoding (fresh support family from the current controlled-transfer state). The no-transfer column is the maximum after-reencoding target-family rate when no basin transfer occurred. Superscripts in the left subtable mark Holm-corrected Wilcoxon significance (\(\ast\), \(p<0.05\)).}
\label{tab:support_diagnostics_refresh}
\begin{subtable}[t]{0.49\textwidth}
  \centering
  \caption{Support diagnostics.}
  \label{tab:fixed17_alignment}
  \label{tab:support_diagnostics}
  \setlength{\tabcolsep}{2pt}
  \resizebox{\linewidth}{!}{\begin{tabular}{@{}l rrrr@{}}
\toprule
Model & $H(B\!\mid\!F_{\rm abs})\downarrow$ & $\overline{|F_{\rm abs}|}$ & wrong supp.\ H1\(\uparrow\) & wrong supp.\ H20\(\uparrow\) \\
\midrule
LISTA & ${\mathbf{0.130}}^{\ast}$ & $\mathbf{4.0}$ & ${2.15{\times}10^{4}}^{\ast}$ & $93.6$ \\
LISTA-BD & ${0.156}^{\ast}$ & $3.8$ & ${2.20{\times}10^{4}}^{\ast}$ & $84.0$ \\
LISTA-SB & ${0.153}^{\ast}$ & $3.8$ & ${2.53{\times}10^{4}}^{\ast}$ & $129$ \\
Sparse MLP, BD & ${0.358}^{\ast}$ & $3.0$ & ${3.87{\times}10^{4}}^{\ast}$ & $93.5$ \\
Sparse MLP & ${0.223}^{\ast}$ & $3.3$ & ${\mathbf{5.82{\times}10^{4}}}^{\ast}$ & ${\mathbf{388}}^{\ast}$ \\
Dense MLP & $1.28$\,\emph{[baseline]} & $1.0$\,\emph{[baseline]} & $17.7$\,\emph{[baseline]} & $2.29$\,\emph{[baseline]} \\
\bottomrule
\end{tabular}
}
\end{subtable}
\hfill
\begin{subtable}[t]{0.49\textwidth}
  \centering
  \caption{Periodic support refresh.}
  \label{tab:support_refresh}
  \setlength{\tabcolsep}{2pt}
  \resizebox{\linewidth}{!}{\begin{tabular}{@{}lccccc@{}}
\toprule
& \multicolumn{4}{c}{Target \(F_{\rm top8}\): stale \(\to\) re-encoded} & No-transfer \\
\cmidrule(lr){2-5}
Model & \(m=1\) & \(m=5\) & \(m=10\) & \(m=20\) & max \(\downarrow\) \\
\midrule
LISTA & $0.999{\to}0.999$ & $0.997{\to}0.999$ & $0.994{\to}1.000$ & $0.963{\to}1.000$ & $0.011$ \\
LISTA-BD & $0.999{\to}0.999$ & $0.995{\to}0.998$ & $0.991{\to}0.999$ & $0.873{\to}1.000$ & $0.013$ \\
LISTA-SB & $0.998{\to}0.999$ & $0.996{\to}0.998$ & $0.989{\to}0.999$ & $0.865{\to}1.000$ & $0.009$ \\
Sparse MLP, BD & $0.996{\to}0.997$ & $0.990{\to}0.994$ & $0.988{\to}0.998$ & $0.833{\to}0.999$ & $0.011$ \\
Sparse MLP & $0.996{\to}0.997$ & $0.993{\to}0.994$ & $0.991{\to}0.998$ & $0.835{\to}0.999$ & $0.015$ \\
Dense MLP & $0.999{\to}1.000$ & $0.977{\to}0.999$ & $0.885{\to}1.000$ & $0.675{\to}1.000$ & $0.211$ \\
\bottomrule
\end{tabular}
}
\end{subtable}
\end{table}

\clearpage
\paragraph{\(F_{\rm abs}\) support families identify basin membership (\cref{tab:support_diagnostics}).}
\begin{figure}[tbp!]
\centering
\includegraphics[width=0.44\linewidth,trim=0 0 392bp 0,clip]{neurips_paper_2026/fig_fixed17_entropy_strips_alt2.pdf}
\caption{Per-seed \(H(B\mid F_{\rm abs})\) on the per-basin deep slice.}
\label{fig:fixed15_entropy_strips}
\label{fig:basin_support_alignment_summary}
\end{figure}
LISTA \(F_{\rm abs}\) support families have low basin uncertainty on the deep-state slice, and sparse-latent MLP models also have far less basin uncertainty than the dense-latent MLP baseline. LISTA rows expose support-family counts close to the benchmark's represented basin count, whereas the dense-latent MLP collapses to a single support family and mixes basins. Sparse-latent encoders therefore expose compact \(F_{\rm abs}\) families that identify basin membership. The per-(system, seed) entropy distribution in \cref{fig:fixed15_entropy_strips} shows that the gap is not driven by outliers.


\paragraph{Accurate forecasts require the correct active coordinates for a basin (\cref{tab:support_diagnostics}).}
The wrong-support ratios are orders of magnitude larger for sparse-latent models than for the dense-latent baseline at horizon 1, and the effect persists at horizon 20.
% Per-system testing rejects on every testable system at the one-step horizon and on the majority at twenty steps; the result is statistically significant.
We conclude that active coordinates are functionally coupled to prediction.

\begin{figure}[H]
\centering
\begin{subfigure}[t]{0.47\linewidth}
  \vspace{0pt}
  \centering
  \includegraphics[width=\linewidth]{neurips_paper_2026/fig_support_coordinate_dropping_accumulated_mse.pdf}
  \caption{Dropping the largest-magnitude active coefficients.}
  \label{fig:support_coordinate_dropping}
\end{subfigure}
\hfill
\begin{subfigure}[t]{0.47\linewidth}
  \vspace{0pt}
  \centering
  \includegraphics[width=\linewidth]{neurips_paper_2026/fig_support_coordinate_random_shuffle_accumulated_mse.pdf}
  \caption{Moving active values to inactive coordinates.}
  \label{fig:support_coordinate_random}
\end{subfigure}

\caption{Inference-time support-coordinate interventions on a representative multibasin system. \textbf{(a)} The standard rollout is compared with cumulative dropping of the largest-magnitude active coefficients of the initial sparse code; shaded bands are bootstrap 95\% intervals over \(100\) starts. \textbf{(b)} Active coefficient values are preserved but reassigned to randomly chosen inactive coordinates; the band is the interquartile range over point-shuffle outcomes.}
\label{fig:support_coordinate_interventions}
\end{figure}

\begin{wrapfigure}[34]{r}{0.44\linewidth}
\vspace{-0.4em}
\centering
\includegraphics[width=\linewidth]{neurips_paper_2026/fig_support_coordinate_trajectories_random_support_19.pdf}
\caption{Wrong-support trajectories. The plot window follows the true and standard rollouts, so divergent intervention paths are clipped.}
\label{fig:support_coordinate_trajectories}
\end{wrapfigure}
The coordinate interventions in \cref{fig:support_coordinate_interventions,fig:support_coordinate_trajectories} test a sharper hypothesis than the cross-system wrong-support mask swap: if the support itself carries basin-relevant dynamical information, then changing only the initial active set should damage forecasts even when the coefficient values are otherwise kept fixed. This is exactly what we observe. At \(H=21\), mean accumulated MSE rises from \(0.0158\) for the standard rollout to \(0.508/1.37/2.08/3.26/8.51\) after zeroing the top \(1/2/3/5/10\) active coefficients ranked by \(|z_{0,i}|\). Random support shuffling is much more destructive: preserving active coefficient values but forcing them onto inactive coordinates yields accumulated errors far above the standard rollout. The trajectory figure shows the qualitative consequence of the same intervention: the standard rollout follows the local basin geometry, whereas wrong-support rollouts move through state-space regions inconsistent with the true basin of attraction. The inference is therefore not merely that sparse codes correlate with basin labels; the active coordinate identities are used by the learned predictor to stay on the correct dynamical branch.

\paragraph{Refreshing \(F_{\rm top8}\) with periodic re-encoding tracks basin transfer (\cref{tab:support_refresh}).}
After controlled basin transfer, re-encoding maps the current state to the target-basin support family. For sparse-family rows, after-reencoding target-family rates are \(0.994\)--\(1.000\) across periods \(1/5/10/20\), while no-transfer false-target rates remain low (\(0.009\)--\(0.019\)). The before/after entries show the mechanism directly: at short periods the stale propagated family is often already target-like, whereas by period \(20\) the stale family degrades to \(0.833\)--\(0.963\) for sparse rows and re-encoding restores it to \(\approx 1.0\). This supports periodic re-encoding as a support-family refresh mechanism without invoking support-conditioned rollouts.

\clearpage
\begin{figure}[tbp!]
\centering
\begin{subfigure}[t]{0.45\linewidth}
  \vspace{0pt}
  \centering
  \includegraphics[width=\linewidth]{neurips_paper_2026/fig_fixed17_horizon_curves.pdf}
  \caption{15-system multibasin benchmark.}
  \label{fig:fixed15_horizon_curves}
  \label{fig:fixed17_horizon_curves}
\end{subfigure}
\hfill
\begin{subfigure}[t]{0.45\linewidth}
  \vspace{0pt}
  \centering
  \includegraphics[width=\linewidth]{neurips_paper_2026/fig_dysts_dt30_forecasting_performance_no_lista_sb.pdf}
  \caption{10-system Dysts $dt{\times}30$ benchmark.}
  \label{fig:dysts_long_horizon}
\end{subfigure}
\caption{Forecasting horizon trends. Lines show arithmetic means across per-system seed-IQM MSEs on a log scale. Bands show fixed-system seed-bootstrap 95\% intervals after system-wise log-relative normalization, anchored to the displayed mean, so they summarize seed uncertainty within an average system.}
\label{fig:forecasting_horizon_curves}
\end{figure}

\section{Conclusion}
\label{sec:conclusion}
Multibasin dynamics make finite-dimensional Koopman learning difficult because the learned representation can hide which basin-scale regime is active. This paper shows that sparse-latent Koopman autoencoders can expose that structure as an interpretable support variable learned without basin supervision. Across controlled multibasin systems, support families align with held-out basin labels, wrong-support interventions show that supports are functionally tied to prediction, and periodic re-encoding refreshes support-family identity after basin transfer. Sparse-latent models also retain forecasting advantages on an external chaotic-flow stress test.

The central takeaway is that sparsity is not only a regularizer for Koopman autoencoders. It can turn a continuous latent representation into a hybrid object: coefficients carry state information, while supports provide discrete, model-produced candidates for basin identity. This suggests a practical path for interpretable multiregime Koopman learning when regime structure is not specified in advance. Future work should scale the support-family mechanism to higher-dimensional systems, learn explicit support-transition graphs, and test whether support variables can be used directly in planning and control.

\appendix

\section{Additional experimental details}
\label{app:experimental_details}

This appendix records the paper-facing experimental protocol. The main experiments ask whether sparse Koopman autoencoders learn useful latent supports without basin supervision. The model is never given basin labels, basin counts, or trajectory-to-basin assignments when training the encoder, decoder, latent transition, support rules, or periodic rollouts. Benchmark basin metadata is used only after training to define evaluation slices, construct controlled diagnostic interventions, and score support--basin agreement. The only training-time exception is architectural: the block-diagonal and soft-block transition diagnostics use benchmark metadata to set the fixed number of transition groups, but they still do not receive trajectory labels or support labels.

\paragraph{Benchmarks.}
The controlled benchmark contains \(15\) two-dimensional multibasin systems, listed with their equations and attractor metadata in \cref{app:fixed17_inventory}. These systems are integrated to produce stored observation sequences; the learner observes only stored states, not vector fields, integration substeps, basin labels, or basin assignments. The external forecasting stress test uses the \(10\) three-dimensional Dysts \(dt{\times}30\) systems listed in \cref{app:dysts_inventory_full}. For Dysts, stored observations are generated at \(30\) times each system's native integration interval and coordinates are standardized after trajectory generation.

\paragraph{Model rows.}
The six model rows in \cref{tab:forecasting_main,tab:fixed17_alignment} are LISTA, LISTA--BD, LISTA--SB, Sparse MLP, Sparse MLP--BD, and Dense MLP. All use latent dimension \(d_z=256\). LISTA rows use shrinkage to produce exact zeros; MLP sparse rows use the sparsity penalty in \cref{eq:loss_latent_terms}; Dense MLP removes the intended sparsity mechanism and sets the sparsity coefficient to zero. The BD rows use a block-diagonal latent transition. The SB row keeps a dense transition but adds an off-block penalty. All rows use a linear decoder dictionary as the state decoder. The controlled multibasin LISTA rows use threshold scale \(\alpha=0.15\), two shrinkage-refinement loops, and a sign-split final operation. The Dysts \(dt{\times}30\) LISTA rows use the same threshold scale with one refinement loop and a ReLU final operation. MLP encoders use two hidden layers of width \(64\); Dense MLP uses a tanh encoder without the sparse output gate.

\paragraph{Controlled multibasin training.}
For every controlled system, each model row is trained for seeds \(0,\ldots,14\). Training uses rollout length \(L=8\) windows, meaning \(9\) adjacent stored states per window, for \(200{,}000\) optimization steps with minibatches of \(256\) windows. The optimizer is AdamW. Encoder and decoder parameters use learning rate \(5{\times}10^{-5}\), the latent transition uses learning rate \(5{\times}10^{-6}\), and non-transition parameters use weight decay \(10^{-4}\). In \cref{eq:training_objective}, the loss weights are \(\lambda_{\rm pred}=1\), \(\lambda_{\rm rec}=0.03\), \(\lambda_{\rm lin}=1\), and \(\lambda_{\rm sp}=3{\times}10^{-3}\) for sparse rows; Dense MLP uses \(\lambda_{\rm sp}=0\). Soft-block rows add the off-block \(L_1\) transition penalty with weight \(10^{-4}\).

All controlled rows use the same label-free boundary-emphasized reset distribution. Before training, \(4096\) candidate initial states are probed for \(32\) steps. Candidates are scored by short-horizon perturbation sensitivity and residual late-rollout motion, and the top \(1024\) states are retained. The scoring probe uses four perturbations, perturbation scale \(0.04\), a late-transient window of \(8\) steps, and transient weight \(0.5\). Half of training windows start from this retained pool with jitter scale \(0.25\). The retained pool is intended to expose boundary-adjacent and transient regions that would otherwise be rare under ordinary resets; it is built without basin labels.

\paragraph{Dysts \(dt{\times}30\) training.}
Dysts rows are trained for the same seeds, \(0,\ldots,14\), with latent dimension \(d_z=256\), minibatches of \(256\), rollout length \(L=10\) windows, and \(100{,}000\) optimization steps. Dysts trajectories come from deterministic native caches with \(200\) cached trajectories, \(30{,}000\) stored observations per trajectory, and a \(2{,}000\)-step warm-up. The first cached trajectory starts from the Dysts default initial condition; remaining trajectories perturb that state with Gaussian noise at scale \(0.2\) times the coordinate-wise Dysts standard deviation. Train, validation, and test caches use separate deterministic namespaces, so held-out test windows are disjoint from training windows and reusable across model seeds. LISTA rows use the controlled-benchmark learning rates. MLP rows use learning rates \(10^{-4}\) for encoder/decoder parameters and \(10^{-5}\) for the latent transition. Sparse Dysts rows use \(\lambda_{\rm sp}=6{\times}10^{-3}\), while Dense MLP uses \(\lambda_{\rm sp}=0\).

\paragraph{Checkpoint selection.}
Checkpoint selection is fixed before test evaluation. Every \(500\) optimization steps, and again at the final step, the current model is rolled out for \(200\) stored steps from \(16\) fixed held-out initial states using every-step re-encoding. The checkpoint with the lowest final validation error under this proxy is used for reported evaluations. Final tables aggregate completed seeds; they do not select the best seed.

\paragraph{Training recipe selection.}
Exploratory sweeps varied LISTA depth, shrinkage scale, sparsity coefficient, sign-split encodings, transition structure, and MLP sparsity controls. These sweeps used shorter budgets or broader system lists and are not pooled with the final evidence. The paper-facing protocol is the fixed recipe in this appendix: \(15\) seeds per retained system, the training budgets above, held-out evaluation splits, and the statistical units in \cref{app:statistical_testing}.

\paragraph{Forecasting experiments.}
Forecasting is evaluated on held-out trajectories at the stored observation cadence. A horizon \(h\) always means \(h\) stored observation steps, not a common physical time across systems. The no-reencoding rollout applies the learned latent transition for the full horizon before decoding. The periodic rollout advances the model's own predicted latent state, decodes the predicted state at a fixed period \(m\), and re-encodes that predicted state; it never uses the true future state. Controlled multibasin forecasting uses \(100\) held-out initial conditions for each system and seed and reports the best score over the fixed period grid \(m\in\{10,25,50,100\}\). Dysts long-horizon forecasting uses \(100\) held-out test-cache windows for each system and seed and reports the best score over \(m\in\{10,25,50,100,150,200\}\). These grids are fixed before aggregation and are not tuned separately for each test trajectory.

\paragraph{Support--basin alignment.}
Support alignment is computed after training from held-out controlled-system states. For each state \(x\), the basin-depth margin is the distance to the second-nearest benchmark attractor center minus the distance to the nearest attractor center. Within each represented benchmark basin, the evaluator keeps the top quartile of held-out states by this margin. This per-basin deep slice uses ground-truth geometry only to define the evaluation population. It is the clean slice for testing basin-support alignment because boundary states can legitimately have unstable or mixed supports.

On this slice, the evaluator encodes states and constructs \(S_{\rm abs}(x)=\{i:|z_i|>10^{-3}\}\). Exact Boolean masks are grouped into \(F_{\rm abs}\) support families by the deterministic greedy Jaccard rule in \cref{sec:method_supports,app:support_definitions} with threshold \(0.5\). The main directional metric is \(H(B\mid F_{\rm abs})\), where \(B\) is the withheld benchmark basin label; lower values mean that knowing the learned support family leaves less uncertainty about basin identity. The companion count \(\overline{|F_{\rm abs}|}\) is descriptive and is interpreted relative to the number of represented basins, not as a monotone score to maximize.

\paragraph{Wrong-support functional diagnostic.}
The cross-system wrong-support diagnostic asks whether the selected active coordinates matter for prediction. For each testable controlled system and basin, the evaluator constructs a canonical deep-slice \(S_{\rm abs}\) mask from the most frequent exact mask in that basin. During a held-out rollout, the inferred mask is replaced with a canonical mask from a different basin and held fixed during the masked latent rollout. \Cref{tab:fixed17_alignment} reports wrong-support/base error ratios at \(H=1\) and \(H=20\); larger ratios mean that using a basin-inconsistent support is more damaging.

\paragraph{Support-coordinate intervention figure.}
The representative coordinate-intervention experiment isolates active-coordinate identity from coefficient values. Starting from a held-out basin-interior state, the model encodes \(z_0=\Enc(x_0)\), perturbs only \(z_0\), and then runs the ordinary learned transition and decoder without retraining. In the coordinate-dropping condition, the active entries are ranked by \(|z_{0,i}|\) and the top \(k\in\{1,2,3,5,10\}\) are set to zero before rollout. In the random-support condition, active coefficient values are preserved but reassigned to randomly chosen inactive coordinates. The figure and table use \(100\) held-out basin-interior starts; the random-support condition uses \(20\) random inactive-coordinate reassignments for each start. Errors are accumulated through \(H=21\), with the full horizon curves shown in \cref{fig:support_coordinate_interventions} and the \(H=21\) summary in \cref{tab:support_coordinate_interventions_h21}.

\paragraph{Periodic support-refresh diagnostic.}
The support-refresh experiment asks whether periodic re-encoding can re-select the support family appropriate to a new region of state space after a controlled transfer. The support object is \(F_{\rm top8}\), built by taking the eight largest-magnitude latent coordinates and merging exact top-\(8\) masks with the same Jaccard-family rule. Each valid comparison starts from a controlled source-to-target transfer generated from benchmark annotations. Between scheduled refresh events, the latent state is advanced only with the learned global transition \(K\). At each event, the stale family implied by the propagated latent is compared with the family obtained by encoding the current controlled-transfer state. The reported cells in \cref{tab:support_refresh} are target-family fractions before re-encoding \(\to\) after re-encoding for fixed periods \(m\in\{1,5,10,20\}\), plus a no-transfer false-target control. Basin labels are used to construct and score transfer instances, not to train a support-conditioned rollout.

\paragraph{Post-hoc routing diagnostic.}
The appendix support-conditioned predictor experiment is a second-stage diagnostic, not part of the main training procedure. After a Koopman autoencoder checkpoint is fixed, the encoder, decoder, and original global \(K\) are frozen. Routes are built from \(S_{\rm top8}\) masks and merged into \(F_{\rm top8}\) support families with Jaccard threshold \(0.40\). Local transition maps are fit by support family on a disjoint fitting split of \(256\) held-out trajectories of length \(256\), giving \(256\times256=65{,}536\) adjacent latent transitions per checkpoint and support definition. A family with fewer than \(50\) current-state fitting transitions does not receive a trainable local map. For each retained family \(c\), the centered affine map \(T_c(z)=\bar z_c+K_c(z-\bar z_c)\) is initialized from the checkpoint's global \(K\) and optimized with decoded rollout MSE while the original autoencoder remains frozen. Route-balanced minibatches prevent the most frequent support families from monopolizing the second-stage updates. Evaluation uses a separate set of \(128\) held-out rollouts; at \(H=1000\), this contributes \(128{,}000\) one-step target positions along evaluated rollouts. Routes are recomputed from the current predicted latent state during rollout, and the routed rollout periodically decodes and re-encodes its own predicted state with period \(5\). If no fitted local route is available, the evaluator falls back to the same model's global transition. The reported routed/global ratio therefore compares a support-routed rollout with the same frozen checkpoint's own global rollout.

\paragraph{Aggregation and significance.}
Point estimates in forecasting and support tables summarize seeds within each system by interquartile mean and then average system summaries across systems. Multibasin forecasting, support-alignment, and wrong-support significance annotations use paired seed-level effects within each controlled system against Dense MLP, with one-sided alternatives fixed by the table direction and Holm correction across eligible systems. Dysts forecasting uses each Dysts system as the independent unit after within-system seed-IQM summarization and applies a paired one-sided exact sign test with Holm correction. Support-refresh and coordinate-intervention panels are descriptive mechanism diagnostics. Full statistical definitions and denominator conventions are in \cref{app:statistical_testing}.


\paragraph{Hardware and resource details.}
Training was done on a SLURM cluster; RTX8000 was the main GPU being used.
Training jobs load CUDA \(12.6.0\) and request one cluster GPU, \(4\) CPU cores,
and \(16\) GB host memory per active training allocation.  Controlled
multibasin training uses single-run GPU allocations with a \(24\)-hour time
limit.  Dysts \(dt{\times}30\) training uses packed GPU allocations with up to
\(12\) runs per allocation, one GPU, \(4\) CPU cores, \(16\) GB memory, and a
\(72\)-hour time limit per packed allocation.  Dysts cache construction uses
CPU allocations with \(4\) CPU cores, \(16\) GB memory, and a \(24\)-hour time
limit.  Dysts long-horizon evaluation uses CPU allocations with \(4\) CPU
cores, \(24\) GB memory, and a \(6\)-hour time limit for the paper-facing
\(dt{\times}30\) evaluation queue.  Controlled support diagnostics and
support-routed analyses use CPU allocations with \(4\) CPU cores, \(16\)--\(24\)
GB memory, and \(8\)--\(12\)-hour time limits, depending on the diagnostic.
\section{Training-dynamics diagnostic}
\label{app:training_dynamics}

The main tables report held-out aggregate performance from validation-selected checkpoints. As a diagnostic for the optimization trajectory, \cref{fig:app_training_dynamics} shows the metric histories for the six controlled multibasin model rows on one representative benchmark system, \texttt{gated\_local\_linear}, with seed \(0\). The curves use the same training protocol and validation rule described in \cref{app:experimental_details}: validation rollouts are evaluated every \(500\) optimization steps, and the saved checkpoint is the step with the lowest validation final error. The diagnostic is not used for model ranking; it is intended to show how the reported checkpoints arise along training.

The validation traces generally improve through training for the sparse rows, while the Dense MLP remains separated by a substantially larger validation error. The best validation checkpoints for the sparse rows occur late in training, between \(179{,}500\) and \(193{,}500\) steps in this trace, whereas the Dense MLP reaches its best validation value earlier. The spectral-radius panel shows that the learned transition remains close to the unit circle throughout training. This supports the qualitative reading that long-horizon behavior is tied to both validation error and the stability of the learned Koopman transition, while avoiding the stronger claim that \(\rho(K)\) closest to one is always optimal.

\begin{figure}[tbp]
\centering
\includegraphics[width=\linewidth]{neurips_paper_2026/appfig_training_dynamics_gated_local_seed0.pdf}
\caption{Representative training-dynamics diagnostic on \texttt{gated\_local\_linear}, seed \(0\), for the six controlled multibasin model rows. Validation final error is the held-out rollout error used by the checkpoint-selection rule; dots mark the minimum validation final error for each trace. The training objective, validation error, transition spectral radius \(\rho(K)\), and sparsity-ratio diagnostic are plotted against optimization step. The horizontal dashed line marks \(\rho(K)=1\). These curves are a diagnostic of training progression and checkpoint selection, not an aggregate benchmark comparison.}
\label{fig:app_training_dynamics}
\end{figure}

For the external chaotic-flow benchmark, \cref{fig:app_training_dynamics_dysts} gives the same diagnostic in aggregate form. Each curve uses seed \(0\) and summarizes the retained 10 Dysts \(dt{\times}30\) systems by the median over finite per-system metric values at each logged step. We use a system median here rather than a single representative Dysts system because the individual chaotic flows have different scales and occasional very large finite validation rollout errors early in training. The resulting validation traces are therefore visibly noisier than the controlled-system trace, but they still show the checkpoint-selection path and the late-training behavior of the model families.

The Dysts diagnostic should be read more cautiously than the controlled benchmark diagnostic. The validation proxy can spike when a chaotic rollout leaves the attractor or compounds error early in training, and Dysts does not provide the controlled basin-label structure used by the main support-alignment claims. Still, the median traces are useful for checking that the retained \(dt{\times}30\) runs are not selected from isolated early accidents: validation curves generally settle after the initial transient, and the learned transition spectra remain close to the unit circle by the end of training, with final median \(\rho(K)\) values roughly between \(1.0003\) and \(1.0011\) across rows. As above, this supports the qualitative stability diagnostic without making \(\rho(K)\) closest to one a sufficient criterion for forecasting quality.

\begin{figure}[tbp]
\centering
\includegraphics[width=\linewidth]{neurips_paper_2026/appfig_training_dynamics_dysts_dt30_seed0.pdf}
\caption{Dysts \(dt{\times}30\) training-dynamics diagnostic. Curves summarize seed \(0\) over the retained 10 Dysts systems by the median over finite per-system metric values at each logged step. Validation final error is the same held-out rollout proxy used by checkpoint selection; dots mark the minimum median validation final error for each trace. The Dysts validation panel is noisier than the controlled trace because chaotic rollout errors can become very large early in training, so this figure is an appendix diagnostic of optimization and selection behavior rather than a main benchmark ranking.}
\label{fig:app_training_dynamics_dysts}
\end{figure}

\clearpage

\section{Multibasin benchmark inventory and system definitions}
\label{app:fixed17_inventory}

\Cref{tab:fixed17_inventory} lists the \(15\) two-dimensional multibasin systems used in this paper. 
Every row contributes to the retained-system aggregate in
\cref{tab:fixed17_alignment}, \cref{tab:self_routing}, and
\cref{fig:fixed17_horizon_curves}. Four systems are additionally displayed as
the main benchmark visual panels in \cref{fig:benchmarks}. The support-refresh
subtable uses \(12\) retained systems for which the controlled-transfer
construction produced eligible post-entry comparisons. 

Basin labels are used
only for benchmark annotation, controlled-transfer construction, and evaluation.
Basin counts are benchmark metadata; for the
structured-transition ablation rows, they are also used only to set a fixed
architecture group count, as described in \cref{app:experimental_details}, and
not for support extraction, support-family construction, routing, or rollout
evaluation.

The 15 multibasin benchmark systems were generated procedurally with Claude Opus 4.5. 

\begin{table}[tbp]
\centering
\caption{The \(15\) multibasin benchmark systems.}
\label{tab:fixed17_inventory}
\resizebox{\textwidth}{!}{%
\begin{tabular}{@{}l l l c@{}}
\toprule
Display name & System key & Generator family & Basins \\
\midrule
Local-linear gates
  & \texttt{gated\_local\_linear}
  & piecewise local-linear gates & 3 \\
Transfer-gated local-linear
  & \texttt{gated\_transfer\_linear}
  & piecewise transfer gates & 3 \\
Arrested spiral
  & \texttt{arrested\_spiral}
  & spiral capture with Gaussian traps & 5 \\
Asymmetric three-well
  & \texttt{cal\_asymmetric\_3}
  & Gaussian wells with independent rotation & 3 \\
High-cross three-well
  & \texttt{cal\_high\_cross\_3}
  & Gaussian wells with stronger rotation & 3 \\
Hexagonal six-well
  & \texttt{cal\_hexagon\_6}
  & Gaussian wells with independent rotation & 6 \\
Octagonal eight-well
  & \texttt{cal\_octagon\_8}
  & Gaussian wells with independent rotation & 8 \\
Pentagonal five-well
  & \texttt{cal\_pentagon\_5}
  & Gaussian wells with independent rotation & 5 \\
Square four-well
  & \texttt{cal\_square\_4}
  & Gaussian wells with independent rotation & 4 \\
Triple-well Duffing
  & \texttt{duffing\_triple\_well}
  & triple-well Duffing oscillator & 3 \\
SNIC multi-attractor
  & \texttt{snic\_multi}
  & polar SNIC-like phase dynamics & 3 \\
Transition-routes four-well
  & \texttt{transition\_routes\_4}
  & route-augmented Gaussian wells & 4 \\
Depth-gradient four-well
  & \texttt{var\_depth\_gradient\_4}
  & Gaussian wells with depth gradient & 4 \\
Diamond four-well
  & \texttt{var\_diamond\_4}
  & Gaussian wells on a diamond layout & 4 \\
L-shaped five-well
  & \texttt{var\_l\_shape\_5}
  & Gaussian wells on an L-shaped layout & 5 \\
\bottomrule
\end{tabular}
}
\end{table}

\paragraph{Stored-step convention.}
All retained systems are implemented as continuous-time vector fields
\(\dot x=f_s(x)\) and integrated with fourth-order Runge--Kutta to produce the
stored observation map \(F_{\Delta t}\) used in \cref{eq:sampled_map}. The
training windows contain only stored states \(x_k\), not \(f_s\), integration
substeps, or basin assignments.

\paragraph{Gaussian-well family.}
Most retained catalog systems use a common two-dimensional Gaussian potential
with an independent rotational drift. For \(x=(x_1,x_2)\in\R^2\), let
\(R x=(x_2,-x_1)\) and
\[
  V_s(x)
  =
  \sum_{i=1}^{M_s}
  -a_{s,i}
  \exp\!\left(-\frac{\norm{x-c_{s,i}}_2^2}{2\sigma_{s,i}^2}\right)
  +\gamma_s(x_1^4+x_2^4).
\]
The implemented vector field is
\begin{equation}
  \dot x = -\nabla V_s(x)+\omega_s R x .
  \label{eq:app_gaussian_well}
\end{equation}
Define \(p_i^{(M)}(r,\varphi)
=r(\cos(2\pi i/M+\varphi),\sin(2\pi i/M+\varphi))\), \(i=0,\ldots,M-1\).
The retained Gaussian-well systems instantiate \cref{eq:app_gaussian_well} with
the parameters in \cref{tab:app_gaussian_params}.

\begin{table}[tbp]
\centering
\caption{Parameters for retained Gaussian-well systems. Repeated pairs in the
\((a_i,\sigma_i)\) column apply to every listed center.}
\label{tab:app_gaussian_params}
\resizebox{\textwidth}{!}{%
\begin{tabular}{@{}l l l l@{}}
\toprule
System & \(\{c_i\}\) & \((a_i,\sigma_i)\) & \((\omega,\gamma)\) \\
\midrule
High-cross three-well
  & \(\{p_i^{(3)}(1.8,\pi/2)\}\)
  & \((3.0,0.5)\) & \((2.0,0.03)\) \\
Hexagonal six-well
  & \(\{p_i^{(6)}(1.7,0)\}\)
  & \((2.0,0.5)\) & \((1.0,0.03)\) \\
Octagonal eight-well
  & \(\{p_i^{(8)}(2.2,0)\}\)
  & \((3.0,0.5)\) & \((0.90,0.02)\) \\
Pentagonal five-well
  & \(\{p_i^{(5)}(1.8,\pi/2)\}\)
  & \((2.0,0.5)\) & \((1.1,0.03)\) \\
Square four-well
  & \(\{p_i^{(4)}(1.8,\pi/4)\}\)
  & \((3.0,0.5)\) & \((1.0,0.03)\) \\
Diamond four-well
  & \(\{(0,2.2),(2.2,0),(0,-2.2),(-2.2,0)\}\)
  & \((2.5,0.5)\) & \((1.0,0.02)\) \\
L-shaped five-well
  & \(\{(-1.5,1.5),(-1.5,0),(-1.5,-1.5),(0,-1.5),(1.5,-1.5)\}\)
  & \((2.5,0.5)\) & \((1.0,0.03)\) \\
Asymmetric three-well
  & \(\{(0,1.8),(-1.5,-0.9),(1.5,-0.9)\}\)
  & \(\{(2.5,0.55),(1.5,0.4),(2.0,0.5)\}\) & \((1.0,0.03)\) \\
Depth-gradient four-well
  & \(\{(-1.3,1.3),(1.3,1.3),(1.3,-1.3),(-1.3,-1.3)\}\)
  & \(\{(2.2,0.55),(2.5,0.5),(3.0,0.5),(3.5,0.5)\}\) & \((1.3,0.03)\) \\
\bottomrule
\end{tabular}
}
\end{table}

The transition-routes system uses the same Gaussian-well form with centers
\[
  \{c_i\}
  =
  \{(-1.8,1.8),(1.8,1.8),(-1.8,-1.8),(1.8,-1.8)\},
\]
\(a_i=3.0\), \(\sigma_i=0.6\), \(\omega=1.0\), and \(\gamma=0.03\), and adds a
corridor-dependent rotation term:
\begin{equation}
  \dot x
  =
  -\nabla V_s(x)
  +
  \bigl(1.0+0.3\,C_{\rm route}(x_2)\bigr) R x,
  \qquad
  C_{\rm route}(x_2)
  =
  e^{-(x_2-1.8)^2/0.3}+e^{-(x_2+1.8)^2/0.3}.
  \label{eq:app_transition_routes}
\end{equation}

\paragraph{Native gated local-linear systems.}
Both native gated systems place three centers on a circular layout, but with
different radii and region definitions. Let
\(Q(\alpha)=\begin{psmallmatrix}\cos\alpha&-\sin\alpha\\
\sin\alpha&\cos\alpha\end{psmallmatrix}\) and
\(\alpha_b=2\pi b/3\), \(b\in\{0,1,2\}\).

For \texttt{gated\_local\_linear}, \(c_b=1.75(\cos\alpha_b,\sin\alpha_b)\).
Inside the radius-\(1.05\) core of basin \(b\),
\[
  \dot x=A_b(x-c_b),\qquad
  A_b=Q(\alpha_b)\widetilde A_b Q(\alpha_b)^\top,
\]
with
\[
\widetilde A_0=\begin{psmallmatrix}-0.9&-1.2\\1.2&-0.9\end{psmallmatrix},
\quad
\widetilde A_1=\begin{psmallmatrix}-1.35&0.2\\-0.3&-0.7\end{psmallmatrix},
\quad
\widetilde A_2=\begin{psmallmatrix}-0.7&-0.1\\0.5&-1.2\end{psmallmatrix}.
\]
Outside the cores, the active sector is the nearest angular sector \(s(x)\) and
the shared gate matrix
\[
  G=\begin{psmallmatrix}-1.35&-0.9\\0.9&-1.35\end{psmallmatrix}
\]
drives the state as \(\dot x=G(x-c_{s(x)})\).

For \texttt{gated\_transfer\_linear}, \(c_b=1.85(\cos\alpha_b,\sin\alpha_b)\).
The implemented radii and widths are
\[
\begin{aligned}
  \rho_{\rm core}&=0.30, &
  \rho_{\rm source}&=0.80, &
  \rho_{\rm exit}&=0.60, &
  \beta_{\rm exit}&=0.72,\\
  w_{\rm chan}&=0.22, &
  \delta_{\rm lane}&=0.28, &
  \rho_{\rm hand}&=0.45 .
\end{aligned}
\]
The core matrices use the same rotated form with templates
\[
\begin{psmallmatrix}-1.0&-1.1\\1.1&-1.0\end{psmallmatrix},\quad
\begin{psmallmatrix}-1.4&0.2\\-0.2&-0.7\end{psmallmatrix},\quad
\begin{psmallmatrix}-0.8&-0.3\\0.5&-1.3\end{psmallmatrix}.
\]
For each ordered pair \(p=(s,t)\), \(s\ne t\), define
\[
  d_{s\to t}=\frac{c_t-c_s}{\norm{c_t-c_s}_2},\qquad
  n_{s\to t}=(-d_{s\to t,2},d_{s\to t,1}),\qquad
  o_t=\frac{c_t}{\norm{c_t}_2},
\]
and \(\chi_{s,t}=1\) when \(\sin(\alpha_s-\alpha_t)\ge 0\), otherwise
\(\chi_{s,t}=-1\). The implemented channel entry and handoff points are
\[
\begin{aligned}
  e_p &= c_s+\rho_{\rm source}d_{s\to t}
        +\chi_{s,t}\delta_{\rm lane}n_{s\to t},\\
  h_p &= c_t+\rho_{\rm hand}o_t
        +\chi_{s,t}\delta_{\rm lane}(-o_{t,2},o_{t,1}),
\end{aligned}
\]
with channel direction \(d_p=(h_p-e_p)/\norm{h_p-e_p}_2\), channel normal
\(n_p=(-d_{p,2},d_{p,1})\), and channel length
\(\ell_p=(h_p-e_p)\cdot d_p\).

The vector field is piecewise analytic. Core regions
\(\norm{x-c_b}_2\le\rho_{\rm core}\) use \(A_b(x-c_b)\). Source annuli use
\(0.65A_b(x-c_b)\), and other non-channel background regions use
\(0.50A_b(x-c_b)\) for the nearest center \(c_b\). Exit wedges for
\(p=(s,t)\) are source-neighborhood states outside the core with
\(\norm{x-c_s}_2\ge\rho_{\rm exit}\) and
\((x-c_s)\cdot d_{s\to t}\ge\norm{x-c_s}_2\cos\beta_{\rm exit}\); they use
\[
  \dot x = v_{\rm exit}d_{s\to t}
  -\lambda_{\rm exit}\bigl((x-c_s)\cdot n_{s\to t}\bigr)n_{s\to t},
  \qquad
  v_{\rm exit}=1.0,\quad \lambda_{\rm exit}=1.8;
\]
and channel rectangles satisfying
\[
  0\le (x-e_p)\cdot d_p\le \ell_p,\qquad
  |(x-e_p)\cdot n_p|\le w_{\rm chan}
\]
use
\[
  \dot x = v_{\rm chan}d_p
  -\lambda_{\rm chan}\bigl((x-e_p)\cdot n_p\bigr)n_p,
  \qquad
  v_{\rm chan}=1.55,\quad \lambda_{\rm chan}=2.8.
\]

\paragraph{Other specialized systems.}
The arrested spiral system has a background spiral flow plus four Gaussian
traps and an origin trap:
\[
  \dot x
  =
  -0.3x+2.0Rx
  -4.0\sum_{i=1}^{4}
  e^{-\norm{x-c_i}_2^2/(2\cdot0.25)}
  \frac{x-c_i}{0.25}
  -1.5e^{-\norm{x}_2^2/0.3}\frac{x}{0.15}
  -0.02(x_1^3,x_2^3),
\]
with \(c_i\in\{(1.5,0),(0,1.8),(-1,-1),(0.8,-1.5)\}\). Its fifth basin is the
origin trap, matching the benchmark manifest count.

The triple-well Duffing system uses \(x=(q,p)\) and
\[
  V(q)=q^6/6-q^4/2+a q^2,\qquad
  V'(q)=q^5-2q^3+2a q,
\]
with \(a=0.3\), damping \(\delta=0.5\), rotation strength \(\omega=1.0\), and
soft cubic confinement \(-\epsilon u^3/s^3\) with \(\epsilon=0.003\) and
\(s=4.0\):
\begin{equation}
\begin{aligned}
  \dot q &= (1+\omega)p-\epsilon q^3/s^3,\\
  \dot p &= -V'(q)-\delta p-\omega q-\epsilon p^3/s^3 .
\end{aligned}
\label{eq:app_duffing_triple}
\end{equation}

The SNIC system is implemented in polar form and then converted to Cartesian
coordinates. With \(\varepsilon_{\rm num}=10^{-8}\),
\(r^2=x_1^2+x_2^2+\varepsilon_{\rm num}\), \(r=\sqrt{r^2}\),
\(\theta=\operatorname{atan2}(x_2,x_1)\),
\(c_\theta=x_1/(r+\varepsilon_{\rm num})\), and
\(s_\theta=x_2/(r+\varepsilon_{\rm num})\),
\[
  \dot r=r(1-r^2)-0.3\cos(3\theta),\qquad
  \dot\theta=1-1.2\cos(3\theta).
\]
The Cartesian vector field is
\[
\begin{aligned}
  \dot x_1
  &= \dot r c_\theta-r\dot\theta s_\theta
     -0.01x_1(x_1^2+x_2^2)+0.5x_2,\\
  \dot x_2
  &= \dot r s_\theta+r\dot\theta c_\theta
     -0.01x_2(x_1^2+x_2^2)-0.5x_1 .
\end{aligned}
\]
These analytic generators define the benchmark trajectories. Held-out
evaluation annotations are produced by the benchmark label helpers or by
endpoint-rollout basin labeling, and the learned models observe only stored
states.

\section{Dysts benchmark inventory and system definitions}
\label{app:dysts_inventory_full}

This appendix lists the ten Dysts systems used in the paper-facing
\(dt{\times}30\) forecasting benchmark. The systems are drawn from Dysts
\citep{gilpin_chaos_2021,gilpin_model_2023}. We used the continuous-time
right-hand sides and metadata from the pinned package version resolved in the
project environment, \texttt{dysts} \(0.96\). All ten systems are
three-dimensional autonomous flows. The equations below are written in
unstandardized native Dysts coordinates \((x,y,z)\); the training cache
standardizes coordinates only after trajectories are generated.

\begin{table}[tbp]
\centering
\caption{The ten Dysts systems used in the \(dt{\times}30\) forecasting
benchmark. The ``stored step'' column is \(30\) times the native Dysts
integration interval. The final column gives the hand-specified diagnostic
structure count used only to size the BD and SB transition parameterizations;
it is a heuristic lobe, scroll, or equilibrium count, not basin truth.}
\label{tab:dysts_inventory_full}
\resizebox{\textwidth}{!}{%
\begin{tabular}{@{}l c c c c@{}}
\toprule
Display name & Dimension & Native \(dt\) & Stored step \(30dt\) & Diagnostic count \\
\midrule
Chua & 3 & \(2.8474745791{\times}10^{-4}\) & \(8.5424237373{\times}10^{-3}\) & 2 \\
Dadras & 3 & \(6.5782963827{\times}10^{-4}\) & \(1.9734889148{\times}10^{-2}\) & 2 \\
Dequan Li & 3 & \(1.6763993999{\times}10^{-5}\) & \(5.0291981997{\times}10^{-4}\) & 3 \\
Hadley & 3 & \(2.9086847808{\times}10^{-4}\) & \(8.7260543424{\times}10^{-3}\) & 3 \\
Lu--Chen--Cheng & 3 & \(1.8469678280{\times}10^{-4}\) & \(5.5409034839{\times}10^{-3}\) & 4 \\
Qi--Chen & 3 & \(7.8371061844{\times}10^{-5}\) & \(2.3511318553{\times}10^{-3}\) & 2 \\
Sakarya & 3 & \(9.9704617436{\times}10^{-4}\) & \(2.9911385231{\times}10^{-2}\) & 2 \\
San--Um--Srisuchinwong & 3 & \(1.4933288881{\times}10^{-3}\) & \(4.4799866644{\times}10^{-2}\) & 2 \\
Shimizu--Morioka & 3 & \(2.4080013336{\times}10^{-3}\) & \(7.2240040007{\times}10^{-2}\) & 2 \\
Wang--Sun & 3 & \(5.3924987498{\times}10^{-3}\) & \(1.6177496249{\times}10^{-1}\) & 4 \\
\bottomrule
\end{tabular}
}
\end{table}

\paragraph{Closed-form convention.}
For each system, the stored observations are generated from \(\dot x=f_s(x)\)
with observation interval \(\Delta t=30\,dt_{\rm native}\), where
\(dt_{\rm native}\) is the system-specific Dysts integration interval in
\cref{tab:dysts_inventory_full}. The learner observes only the resulting stored
states, not the vector field, continuous-time solver, or substeps. All ten
systems have explicit closed-form right-hand sides in Dysts. Most are
polynomial vector fields. The Chua and San--Um--Srisuchinwong systems include
absolute-value terms, so they are piecewise smooth rather than globally
analytic at the corresponding switching surfaces.

\paragraph{Chua.}
The Chua system uses the piecewise-linear diode nonlinearity
\[
  h(x)
  =
  m_1x+\frac{1}{2}(m_0-m_1)\bigl(|x+1|-|x-1|\bigr),
\]
with parameters
\[
  \alpha=15.6,\qquad
  \beta=28,\qquad
  m_0=-1.142857,\qquad
  m_1=-0.71429 .
\]
The vector field is
\begin{equation}
\begin{aligned}
  \dot x &= \alpha\bigl(y-x-h(x)\bigr),\\
  \dot y &= x-y+z,\\
  \dot z &= -\beta y .
\end{aligned}
\label{eq:app_dysts_chua}
\end{equation}

\paragraph{Dadras.}
With parameters
\[
  c=2,\qquad e=9,\qquad o=2.7,\qquad p=3,\qquad r=1.7,
\]
the Dadras system is
\begin{equation}
\begin{aligned}
  \dot x &= y-px+oyz,\\
  \dot y &= ry-xz+z,\\
  \dot z &= cxy-ez .
\end{aligned}
\label{eq:app_dysts_dadras}
\end{equation}

\paragraph{Dequan Li.}
With parameters
\[
  a=40,\qquad c=1.833,\qquad d=0.16,\qquad
  \varepsilon=0.65,\qquad f=20,\qquad k=55,
\]
the Dequan Li system is
\begin{equation}
\begin{aligned}
  \dot x &= a(y-x)+dxz,\\
  \dot y &= kx+fy-xz,\\
  \dot z &= cz+xy-\varepsilon x^2 .
\end{aligned}
\label{eq:app_dysts_dequan_li}
\end{equation}

\paragraph{Hadley.}
With parameters
\[
  a=0.2,\qquad b=4,\qquad f=9,\qquad g=1,
\]
the Hadley system is
\begin{equation}
\begin{aligned}
  \dot x &= -y^2-z^2-ax+af,\\
  \dot y &= xy-bxz-y+g,\\
  \dot z &= bxy+xz-z .
\end{aligned}
\label{eq:app_dysts_hadley}
\end{equation}

\paragraph{Lu--Chen--Cheng.}
With parameters
\[
  a=-10,\qquad b=-4,\qquad c=18.1,
\]
the Lu--Chen--Cheng system is
\begin{equation}
\begin{aligned}
  \dot x &= -\frac{ab}{a+b}x-yz+c,\\
  \dot y &= ay+xz,\\
  \dot z &= bz+xy .
\end{aligned}
\label{eq:app_dysts_lu_chen_cheng}
\end{equation}

\paragraph{Qi--Chen.}
With parameters
\[
  a=38,\qquad b=2.666,\qquad c=80,
\]
the Qi--Chen system is
\begin{equation}
\begin{aligned}
  \dot x &= a(y-x)+yz,\\
  \dot y &= cx+y-xz,\\
  \dot z &= xy-bz .
\end{aligned}
\label{eq:app_dysts_qi_chen}
\end{equation}

\paragraph{Sakarya.}
With parameters
\[
  a=-1,\qquad b=1,\qquad c=1,\qquad h=1,\qquad
  p=1,\qquad q=0.4,\qquad r=0.3,\qquad s=1,
\]
the Sakarya system is
\begin{equation}
\begin{aligned}
  \dot x &= ax+hy+syz,\\
  \dot y &= -by-px+qxz,\\
  \dot z &= cz-rxy .
\end{aligned}
\label{eq:app_dysts_sakarya}
\end{equation}

\paragraph{San--Um--Srisuchinwong.}
With parameter \(a=2\), the San--Um--Srisuchinwong system is
\begin{equation}
\begin{aligned}
  \dot x &= y-x,\\
  \dot y &= -z\tanh(x),\\
  \dot z &= -a+xy+|y| .
\end{aligned}
\label{eq:app_dysts_san_um_srisuchinwong}
\end{equation}

\paragraph{Shimizu--Morioka.}
With parameters
\[
  a=0.85,\qquad b=0.5,
\]
the Shimizu--Morioka system is
\begin{equation}
\begin{aligned}
  \dot x &= y,\\
  \dot y &= x-ay-xz,\\
  \dot z &= -bz+x^2 .
\end{aligned}
\label{eq:app_dysts_shimizu_morioka}
\end{equation}

\paragraph{Wang--Sun.}
With parameters
\[
  a=0.2,\qquad b=-0.01,\qquad d=-0.4,\qquad
  e=-1,\qquad f=-1,\qquad q=1,
\]
the Wang--Sun system is
\begin{equation}
\begin{aligned}
  \dot x &= ax+qyz,\\
  \dot y &= bx+dy-xz,\\
  \dot z &= ez+fxy .
\end{aligned}
\label{eq:app_dysts_wang_sun}
\end{equation}


\section{Support definitions and sensitivity}
\label{app:support_definitions}

\Cref{sec:method_supports} defines the notation used in the main text. This appendix records the implementation-level objects behind that notation and explains how to read the corresponding diagnostics. All support objects are computed after training from the learned encoder output \(z=\Enc(x)\). Basin labels, basin counts, and trajectory-to-basin assignments are not used to train the model, define a support, or build a support family; when labels appear below they are evaluation-only quantities used to score or display the learned objects.

\paragraph{From a latent vector to an exact support key.}
The reducers first convert each latent vector into a Boolean mask \(m(x)\in\{0,1\}^{d_z}\). The exact support \(S(x)\) is the index set of the true entries of this mask, but the code compares exact supports by serializing the Boolean mask itself: two states have the same exact support if and only if all \(d_z\) Boolean entries match. The main support rule used in the paper is
\begin{align}
  m_{{\rm abs},i}(x)
  &=\mathbf{1}\{|z_i|>10^{-3}\},
  &
  S_{\rm abs}(x)
  &=\{i:m_{{\rm abs},i}(x)=1\},
\end{align}
with \(\epsilon=10^{-12}\) in the implementation. The inequalities are strict. The epsilon only affects the degenerate all-zero relative-threshold case, where \(S_{\rm rel}\) is empty. \(S_{\rm abs}\) is a magnitude-thresholded active set, so \(|S_{\rm abs}|\) is a measured state-dependent statistic rather than a fixed sparsity budget. 
% \(S_{\rm top8}\) instead always begins from eight active coordinates when \(d_z\ge 8\); it ignores absolute scale and gives every state a fixed-size exact key before family merging.

\paragraph{From exact supports to support families.}
For a chosen support rule and evaluation collection, the implementation counts distinct exact support keys, sorts them from most to least frequent, and uses the serialized key as a deterministic tie-breaker. It then performs the leader-style greedy Jaccard merge described in \cref{sec:method_supports}. Each family stores an actual Boolean support vector as its representative. A new exact mask joins the existing representative with the largest Jaccard similarity if that similarity is at least \(0.5\); otherwise it starts a new family. Representatives are not optimized, averaged, or updated as centroids. Family labels are therefore run- and collection-specific integer identifiers, and only their induced partition of states is meaningful.

\paragraph{Support-family creation algorithm.}
For completeness, the family construction used for both \(F_{\rm abs}\) and \(F_{\rm top8}\) is:
\begin{quote}
\small
\textbf{Input:} evaluation states \(\mathcal X\), a support rule \(q\), and Jaccard threshold \(\tau=0.5\).\\
\textbf{Output:} family label \(F_q(x)\) for each \(x\in\mathcal X\), and family representatives \(\{r_f\}\).
\begin{enumerate}[label=\textbf{\arabic*.},leftmargin=2.0em]
  \item Compute one exact Boolean mask \(m_q(x)\in\{0,1\}^{d_z}\) for every \(x\in\mathcal X\).
  \item Count the multiplicity \(c(m)\) of each distinct exact mask \(m\).
  \item Sort the distinct masks in decreasing \(c(m)\), using the serialized Boolean mask as a deterministic tie-breaker.
  \item Initialize an empty list of families and representatives.
  \item For each distinct mask \(m\) in the sorted order:
  \begin{enumerate}[label=\textbf{\alph*.},leftmargin=2.0em]
    \item If no family exists, create a new family \(f\), set \(r_f\gets m\), and assign \(m\) to \(f\).
    \item Otherwise compute \(f^\star=\operatorname*{arg\,max}_f J(m,r_f)\) over existing family representatives.
    \item If \(J(m,r_{f^\star})\ge\tau\), assign every occurrence of \(m\) to family \(f^\star\).
    \item If \(J(m,r_{f^\star})<\tau\), create a new family \(f\), set \(r_f\gets m\), and assign every occurrence of \(m\) to \(f\).
  \end{enumerate}
  \item For each state \(x\), return \(F_q(x)\), the family assigned to its exact mask \(m_q(x)\).
\end{enumerate}
\end{quote}
The representative update \(r_f\gets m\) occurs only when a new family is created. Existing representatives are never averaged, recomputed as majority masks, or moved toward later assigned masks.

For top-eight masks, \(J(m,r)\ge 0.5\) is equivalent to at least six shared active coordinates because both masks have size eight. For \(S_{\rm abs}\), the same threshold is adaptive to mask size through the intersection-over-union denominator. A lower Jaccard threshold merges more exact supports and can collapse distinct basins; a higher threshold separates more masks and can make \(H(B\mid F)\) small by over-fragmenting each basin. We therefore fix \(0.5\) rather than tuning it post hoc and interpret any entropy result together with the corresponding family count.

\paragraph{Stable support component construction.}
\(C_{\rm stab}\) is used in the staged local-linearization experiment, not as the primary object in the basin-support alignment table. It also starts from the absolute mask \(m_{\rm abs}(x)\), but the first label is not \(F_{\rm abs}\). For this experiment, exact \(m_{\rm abs}\) masks are merged into a finer base support-family label \(u_t\) with Jaccard threshold \(0.8\). Given a training-split trajectory, \(u_{i,t}\) denotes the base label of trajectory \(i\) at stored time \(t\). A transition \(u\to v\) means that consecutive stored observations satisfy \(u_{i,t}=u\) and \(u_{i,t+1}=v\). The empirical transition graph is built from 512 trajectories of length 192 after the first half of LISTA training, while the encoder and decoder are held fixed.

The graph is filtered to retain sufficiently frequent transitions, its strongly connected components are computed, and recurrent components are selected when they appear in trajectory tails and have small outbound transition mass. Each observed base label \(u\) is then assigned to the recurrent component that its future support trajectory reaches with high empirical confidence. This assignment is \(C_{\rm stab}(u)\). Thus \(C_{\rm stab}\) groups base support families by support-flow fate; it is neither a ground-truth basin label nor an estimate of a state-space fixed point.

\paragraph{Dense tail-fate control.}
The dense control used to audit the local-linearization result deliberately removes the sparse support graph. It starts from the dense zero-sparsity MLP KAE, encodes route-fitting trajectories after the global training stage, and summarizes each dense latent trajectory by the mean, standard deviation, and final value of its tail. These continuous tail summaries are standardized, optionally PCA-compressed, and clustered with \(k\)-means using a silhouette-selected cluster count up to \(12\). All transitions from a trajectory inherit that trajectory's dense fate cluster. During forecasting, the future trajectory tail is unavailable, so the route is assigned by nearest centroid in the current dense latent space, using the mean current latent state for each fitted dense fate label. This is the intended contrast with \(C_{\rm stab}\): the dense control asks which dense fate centroid the current state is closest to, whereas \(C_{\rm stab}\) asks which recurrent support-flow component the current sparse support pattern is assigned to. Both objects can contain basin-fate information, but only \(C_{\rm stab}\) exposes an active-coordinate support graph and routes from the current support mask.

\paragraph{Which support object is used where.}
\begin{description}[leftmargin=1.6em,style=nextline]
  \item[\(S_{\rm abs}\).]
  This is the strict active-coordinate object. It is used to build \(F_{\rm abs}\), to report exact-support diagnostics such as \(H(B\mid S_{\rm abs})\), \(H(S_{\rm abs}\mid B)\), exact-support uniqueness, and \(|S_{\rm abs}|\), and to form canonical masks for wrong-support interventions. In the intervention code, the canonical mask for basin \(b\) is the most frequent \(S_{\rm abs}\) key among candidate deep-slice states from basin \(b\). The ``wrong'' condition replaces it with canonical masks from other represented basins and holds that mask fixed during the masked latent rollout.

  \item[\(F_{\rm abs}\).]
  This is the main basin-support alignment object in \cref{tab:fixed17_alignment}. It asks whether the many exact \(S_{\rm abs}\) masks produced by a sparse encoder compress into basin-scale families. The primary directional metric is \(H(B\mid F_{\rm abs})\): low values mean that knowing the support family leaves little uncertainty about the withheld basin label. The count \(\overline{|F_{\rm abs}|}\) is deliberately non-directional and should be compared with the number of basin labels represented in the evaluated slice.

  \item[\(C_{\rm stab}\).]
  This is the route object for the staged local affine maps in \cref{sec:cstab_local_linearizations}. During the second half of training, a latent source state is converted to \(m_{\rm abs}\), assigned to a base support-family label \(u_t\), mapped to \(c=C_{\rm stab}(u_t)\), and then used to choose the affine latent map \(T_c(z)=d_c+K_c(z-\bar z_c)\). All local maps are trained in parallel with the encoder and decoder frozen. During forecasting with periodic re-encoding, the route is recomputed from the model's current predicted state. Exact support masks observed when \(C_{\rm stab}\) was built map directly to their component; unseen masks are assigned to the nearest component prototype when their Jaccard overlap is at least \(0.4\); otherwise the rollout uses the frozen global \(K\) step.
\end{description}

\paragraph{How to read the diagnostics.}
A low \(H(B\mid S_{\rm abs})\) or \(H(B\mid F_{\rm abs})\) means that the support object predicts basin identity on the evaluated benchmark states. It does not by itself imply a compact one-support-per-basin code: a model can achieve low \(H(B\mid S_{\rm abs})\) while assigning many different exact masks to the same basin. \(H(S_{\rm abs}\mid B)\), exact-support uniqueness, and \(|S_{\rm abs}|\) diagnose this fragmentation, while \(H(B\mid F_{\rm abs})\) and \(\overline{|F_{\rm abs}|}\) ask whether the fragmented exact masks merge into a basin-scale family structure.

The takeaways are as follows. \(F_{\rm abs}\) is the object for the claim that sparse supports identify basin interiors. \(S_{\rm abs}\) is the intervention object for testing whether the active coordinates are functionally used. \(C_{\rm stab}\) is the route object for the staged local-linearization result; it is built from support transitions after stage-one training and then used to select local affine latent dynamics during stage-two training and periodic-reencoding rollout.

\section{Statistical testing protocol}
\label{app:statistical_testing}

Point estimates and significance annotations answer different questions. Forecasting and alignment entropy tables first summarize finite completed seeds within each system by \texttt{scipy.stats.trim\_mean} with proportion \(0.25\), then take an arithmetic mean across systems. This sorts \(n\) values and discards \(\lfloor n/4\rfloor\) from each tail; complete \(15\)-seed cells therefore average their central nine. We refer to this implementation as the interquartile mean (IQM). Forecasting averages all \(15\) controlled systems; primary alignment averages the \(14\) systems with at least two observed labels and retains an all-\(15\) sensitivity. The descriptive family-count column instead averages per-system seed means. We do not pool trajectories, seeds, and systems into one inferential sample.

\paragraph{Forecasting metric and missingness.}
For start \(i\), rollout step \(t\), and state coordinate \(j\), let
\[
  e_{i,t}=\sum_{j=1}^{d_x}(\hat x_{i,t,j}-x_{i,t,j})^2.
\]
For the KAE rows used in formal comparisons, the reported MSE through horizon \(H\) takes the mean of finite \(e_{i,t}\) over \(t=1,\ldots,H\) within each start, then averages finite per-start means. It is therefore a state-summed squared error, not per-coordinate MSE. A rollout can contribute through its finite prefix even if it is nonfinite before \(H\), so the KAE ranking does not enforce full-horizon finiteness. Local-linear mixture controls use the same finite-step convention; the unmatched classical DMD/EDMD evaluator requires a finite ordinary through-\(H\) mean for each contributing start, as detailed in \cref{app:experimental_details}.

The error \(E_{s,r,h}\) used below is also best-periodic: for every system \(s\), seed \(r\), model, and horizon \(h\), the lowest MSE is selected on the same test trajectories over the fixed cadence grid. This test-set oracle is shared by the point estimates and formal KAE comparisons and should not be interpreted as validation-frozen deployment selection.

\paragraph{Controlled forecasting effects.}
For candidate model \(a\), Dense MLP baseline, system \(s\), seed \(r\), and horizon \(h\), the paired effect is
\[
  \Delta_{s,r,h}^{a}
  =
  \log_{10}E^{a}_{s,r,h}
  -
  \log_{10}E^{\rm Dense}_{s,r,h}.
\]
Only mutually finite positive MSEs are eligible. Within each system and candidate--horizon cell, a one-sided paired Wilcoxon signed-rank test uses the common seeds with alternative \(\Delta<0\). At least four finite pairs are required. Raw per-system \(p\)-values are Holm-corrected across eligible systems in that candidate--horizon family. A main-table star denotes corrected \(p<0.05\); every displayed controlled sparse-row forecasting cell passes on all \(15/15\) systems.

The per-system tables in \cref{app:controlled_results} display the arithmetic mean of the paired seed effects \(\Delta_{s,r,h}^{a}\), together with the Wilcoxon/Holm \(p\)-value. The displayed effect is a mean, not the median used by an older forest-plot artifact. Because the signed-rank test depends on ranks whereas the mean does not, an extreme finite outlier can make their directions appear inconsistent; the Transfer-gated local-linear LISTA--BD H500 cell is the one such case, with \(12/15\) negative paired effects but a positive arithmetic mean.

\paragraph{Transductive alignment effects.}
The alignment partition \(F_{\rm abs}^{\rm align}\) is fit on all states in an evaluation trajectory collection and scored on states at or above the within-label \(75\)th center-margin percentile contained in that collection; ties are retained, so the subset need not contain exactly one quarter of a label, and no between-label count balancing is performed. Labels and centers are native for the two gated systems. For the \(13\) catalog systems they are nearest-center proxies obtained by deterministic clustering of long-rolled endpoints at the known benchmark count. The margin \(d_2-d_1\) is nearest-center separation, not a true boundary-distance measurement. \(H(B\mid F_{\rm abs}^{\rm align})\) uses natural logarithms and is reported in nats. The seed is the paired unit within each system and the one-sided alternative is candidate \(<\) Dense MLP. Holm correction is across eligible systems within each model--metric cell. Triple-well Duffing's frozen proxy realizes one observed label, so its zero entropy for every model is uninformative and it is excluded from primary aggregates; the all-zero paired differences and undefined signed-rank statistic are consequences. Every displayed alignment star passes on all \(14/14\) eligible systems. The observed high-margin-slice family count has no directional test because closeness to the number of represented evaluation labels, rather than monotone increase, is the relevant descriptive comparison. These tests quantify model differences in the transductive diagnostic; they do not turn it into an out-of-sample classification experiment.

\paragraph{Dysts forecasting effects.}
The historical Dysts KAE rows and their associated point estimates, sign tests, Holm corrections, and significance markers are quarantined because the KAE and standalone-control trajectories used mismatched physical cadences. No Dysts inferential result is reported until the corrected all-six-row campaign is complete.

\paragraph{Descriptive comparisons.}
The standalone DMD, EDMD, and local-linear controls configure three seeds and use independently generated evaluation trajectories; their per-system summaries omit nonfinite seed values, so they are not included in paired KAE tests. The staged support-family result reports \(225\) system--seed pairs nested within \(15\) systems, plus MSE ratios, without a confirmatory \(p\)-value. Its checkpoint score applies the same finite-prefix state-summed error reducer as the reported table, minimizes separately over the eight cadences at H100, H500, and H1000, and averages those three minima. The \(32\) selector starts are the first \(32\) of the \(100\) reported starts, whereas the global baseline selects its checkpoint on \(16\) distinct starts with an H200 final-error proxy. This \(32\%\) staged overlap and selector asymmetry make the comparison optimistic and preclude interpreting the paired win counts as held-out causal evidence. The Local-linear gates coordinate intervention uses one LISTA checkpoint at seed \(0\); bootstrap intervals summarize its \(100\) starts, while random-support medians and interquartile ranges pool \(2{,}000\) dependent point--shuffle outcomes per horizon. These describe within-checkpoint variability only.

\paragraph{High-dimensional stress tests.}
The Lorenz--96, Allen--Cahn temporal-support, and forecast-optimized
Allen--Cahn packets each average trajectories within model seed and use ten
paired model-initialization seeds as inferential units; trajectories are not
independent replicates. Lorenz--96 and the temporal-support packet use
100,000 percentile bootstrap resamples of seed-level outcomes and exact
two-sided sign-flip tests over all \(2^{10}\) assignments. Lorenz--96 applies
Holm correction over its two declared long physical horizons, \(2.5\) and
\(5\); temporal-support forecasting corrects over physical times
\(8,12,16,20\), and its final-state support gate combines frozen absolute
thresholds with a paired sparse-minus-dense ARI interval. The separate
forecast-optimized Allen--Cahn packet uses 50,000 paired-seed bootstrap
resamples of each ratio-of-arm-means reduction. Those four intervals are
marginal; the frozen primary decision is their four-cell intersection-union
gate. An exact one-sided paired sign-flip max-\(t\) calculation over the four
H160/H200 mean/terminal cells is reported only as a multiplicity sensitivity.
H160 is a prefix of H200, so these cells are dependent. Every interval
conditions on its packet's fixed generated data and does not cover new dataset
draws.

The same-checkpoint Allen--Cahn new-initial-condition check crosses three newly
generated datasets with each of the same ten paired sparse/dense checkpoint
seeds. It first averages the three datasets within checkpoint seed, then uses
the ten paired seeds for a 100,000-resample paired bootstrap interval and exact
one-sided sign-flip test over all $2^{10}$ assignments. Thus 30 seed--dataset
cells are not treated as independent replicates, and the three dataset-specific
effects are descriptive. The H200 through-horizon endpoint was frozen only
after the original four-cell gate failed, so this inference is an outcome-aware
same-checkpoint robustness check rather than independent retraining replication
or a repair of that gate.

The secondary Allen--Cahn physical-metric packet uses the identical forecasts
and the same within-checkpoint averaging over three datasets. Its sole
statistical family contains all seven H200 cumulative phase, interface, and
energy metrics. For each metric, the one-sided paired sign-flip test enumerates
all $2^{10}$ assignments of the oriented seed effects; Holm correction is over
all seven raw values. The 100,000-resample paired-seed interval is for the
oriented absolute arm-mean difference. Relative error reductions divide that
absolute difference by the observed dense arm mean; modal accuracy is instead
reported in percentage points. Instantaneous curves, terminal values, H160,
and the late half are mandatory descriptives and receive no separate
inference. Because the metrics translate an outcome-selected H200 advantage
on unchanged checkpoints, favorable results remain secondary and cannot
reclassify the original gate.

\paragraph{Holm correction and assumptions.}
For each formal family, eligible raw \(p\)-values are sorted, multiplied by the number of remaining hypotheses, made monotone by a running maximum, and clipped at one. Paired Wilcoxon tests assume the common seed differences are the relevant exchangeable units within a controlled system. All conclusions condition on the fixed benchmark roster. Evaluation-set period selection, finite-prefix omission, staged checkpoint-selection overlap, selector asymmetry, and transductive family fitting are protocol limitations independent of the multiple-comparison correction.

\section{Post-hoc support-conditioned predictor diagnostics}
\label{app:support_conditioned_predictors}

This appendix records the post-hoc support-conditioned predictor experiments. These experiments are diagnostic: they ask whether learned support families can select fitted transition maps after the Koopman autoencoder has already been trained. They are not used for the main paper's basin-identifiability claim or for the main forecasting table.

\paragraph{Support-conditioned rollouts.}
Support-conditioned prediction is a second-stage procedure applied after the Koopman autoencoder has been trained with a single global transition \(K\). In the second stage, the encoder \(\Enc\), decoder \(\Dec\), and original \(K\) are frozen. A separate centered transition \(K_c\) is optimized for each retained \(F_{\rm top8}\) support family \(c\), while the original \(K\) remains the fallback for unretained or unassigned routes. The comparison point is the same frozen checkpoint's global rollout,
\begin{equation}
  \hat z_{t+1}^{\rm global}=K\hat z_t,\qquad
  \hat x_{t+1}^{\rm global}=\Dec(\hat z_{t+1}^{\rm global}).
  \label{eq:global_rollout}
\end{equation}

Routes are constructed without basin labels. On disjoint fitting trajectories, we encode states with the frozen encoder, compute exact \(S_{\rm top8}\) supports, and merge them into support families by the Jaccard rule in \cref{sec:method_supports} with threshold \(0.40\). Families with fewer than \(50\) current-state fitting transitions are not given trainable maps. For each retained family \(c\in\mathcal C_{\rm fit}\), the fixed center \(\bar z_c\) is the mean current latent assigned to that family, and the map is initialized from the checkpoint transition, \(K_c\leftarrow K\). The route-local transition is
\begin{equation}
  T_c(z)=\bar z_c+K_c(z-\bar z_c).
  \label{eq:family_local_map}
\end{equation}

The reported local maps are trained jointly as one bank of matrices by gradient descent on decoded rollout MSE. For a second-stage training window \((x_{b,0},\ldots,x_{b,H})\), let \(\hat z_{b,0}=\Enc(x_{b,0})\), where \(H\) is the checkpoint's short training horizon. At rollout offset \(\ell\), the current predicted latent is detached only for the discrete route assignment:
\[
  c_{b,\ell}=F_{\rm top8}(\hat z_{b,\ell}).
\]
If this assigned family is retained, the next latent uses \(T_{c_{b,\ell}}\); otherwise it falls back to the frozen global transition:
\begin{equation}
  \tilde z_{b,\ell+1}=
  \begin{cases}
    T_{c_{b,\ell}}(\hat z_{b,\ell}), & c_{b,\ell}\in\mathcal C_{\rm fit},\\
    K\hat z_{b,\ell}, & c_{b,\ell}\notin\mathcal C_{\rm fit}.
  \end{cases}
  \label{eq:family_local_rollout}
\end{equation}
The decoded prediction is \(\hat x_{b,\ell+1}=\Dec(\tilde z_{b,\ell+1})\). Every \(m\) predicted steps, with \(m=5\) in the selected runs, the next latent state is refreshed by the frozen encoder,
\[
  \hat z_{b,\ell+1}=\Enc(\hat x_{b,\ell+1}),
\]
and otherwise \(\hat z_{b,\ell+1}=\tilde z_{b,\ell+1}\). The active setting recomputes \(c_{b,\ell}\) before every one-step map application, so route switching can occur at any predicted step, not only at refresh boundaries.

The second-stage objective is
\begin{equation}
  \min_{\{K_c:\,c\in\mathcal C_{\rm fit}\}}
  \frac{1}{BH}\sum_{b=1}^{B}\sum_{\ell=1}^{H}
  \norm{\hat x_{b,\ell}-x_{b,\ell}}_2^2 .
  \label{eq:stage2_local_loss}
\end{equation}
There is no basin-label loss, support-classification loss, latent consistency loss, or regularizer in this stage. Minibatches are route-balanced: a training step samples support-family buckets approximately uniformly and then samples windows inside the selected buckets. Thus frequent support families do not monopolize the updates, and a particular \(K_c\) receives gradients only on examples and rollout steps where it is selected.

At evaluation time no parameters are updated. Starting from a held-out initial state, the same recurrence in \cref{eq:family_local_rollout} is run for the forecast horizon using only predicted states. Exact supports seen in the codebook map directly to their family; unseen exact supports are assigned to the nearest family prototype only when the Jaccard similarity reaches the same threshold, and otherwise use the global fallback. Exact \(S_{\rm top8}\)-local routing uses the same rollout structure with exact supports as routes and is retained as an ablation in \cref{app:routing_robustness}.

\begin{table}[tbp]
\centering
\caption{Appendix support-conditioned predictor diagnostic. Entries are label-free \(F_{\rm top8}\)-local routed/global MSE ratios and system-win counts on all states; lower ratios and higher win counts are better. Superscripts denote system-level significance tests: \(\ast\), \(p_{\rm Holm}<0.05\); \(\ast\ast\), \(p_{\rm Holm}<0.01\).}
\label{tab:self_routing}
\setlength{\tabcolsep}{2pt}
\resizebox{0.78\textwidth}{!}{\begin{tabular}{@{}l cc cc@{}}
\toprule
& \multicolumn{2}{c}{$H100$} & \multicolumn{2}{c}{$H1000$} \\
\cmidrule(lr){2-3}\cmidrule(l){4-5}
Model & \(F_{\rm top8}\) ratio & system wins & \(F_{\rm top8}\) ratio & system wins \\
\midrule
LISTA & ${0.239}^{\ast}$ & $13/15$ & ${\mathbf{1.41\times10^{-11}}}^{\ast\ast}$ & $15/15$ \\
LISTA-SB & ${0.021}^{\ast}$ & $13/15$ & ${2.42\times10^{-7}}^{\ast}$ & $14/15$ \\
LISTA-BD & ${\mathbf{0.011}}^{\ast}$ & $13/15$ & ${6.96\times10^{-5}}^{\ast}$ & $15/15$ \\
Sparse MLP, BD & ${0.032}^{\ast\ast}$ & $12/15$ & ${6.53\times10^{-8}}$ & $15/15$ \\
Sparse MLP & ${0.038}^{\ast\ast}$ & $13/15$ & ${5.19\times10^{-4}}$ & $15/15$ \\
Dense MLP & ${154}$ & $1/15$ & ${85.7}$ & $1/15$ \\
\bottomrule
\end{tabular}}
\end{table}

\section{Falsification and partition controls}
\label{app:falsification}

The partition-control experiments are falsification checks for the support-routing interpretation in \cref{tab:self_routing}. They separate four questions. First, do post-hoc route-local centered affine maps \(T_c(z)=\bar z_c+K_c(z-\bar z_c)\) improve over the same model's autonomous global-\(K\) rollout at all? Second, can the model-produced \(F_{\rm top8}\) support family select those maps without basin labels or a known basin count? Third, would an oracle basin partition, a generic latent-space partition, or an arbitrary count-matched partition explain the same effect? Fourth, does improvement over the same model's autonomous global-\(K\) rollout imply that the post-hoc routed predictor is also better than the best-periodic re-encoding forecaster used in the main forecasting tables?

All rows in \cref{tab:self_routing_partition_controls} use the current \(d_z=256\) artifacts and freeze the trained encoder \(\Enc\), decoder \(\Dec\), and learned global transition \(K\). The route-operator fitting split contains \(256\) held-out trajectories with \(256\) adjacent encoded transitions each. For each selector, the current latent states define route classes \(c\); for every class with at least \(128\) fitting transitions, we fit a centered ridge slope \(K_c\), with \(\eta=10^{-4}\), and roll out the affine map \(T_c\) in \cref{eq:family_local_map,eq:family_local_rollout}. Missing or underpopulated routes fall back to the same model's global \(K\). The forecast split is disjoint from the fitting split, contains \(128\) held-out rollouts, and recomputes the route from the model's current predicted latent state during rollout. Thus the table changes only the route selector while holding the trained representation, local-map fitting rule, fallback rule, and global baseline fixed.

The selectors differ only in how they define the route class \(c_t\). The support-family selector uses \(c_t=F_{\rm top8}(z_t)\), constructed by greedy Jaccard merging of exact \(S_{\rm top8}\) masks at threshold \(0.5\), and uses no basin labels or basin count. The oracle-basin selector uses benchmark basin labels on the fitting trajectories and, during rollout, assigns the decoded current prediction \(\Dec(z_t)\) to the nearest evaluation basin center used by the labeling routine. This is an evaluation-only control: it uses information that is intentionally unavailable in the training/deployment setting. The latent-cluster selector fits \(k\)-means clusters in the learned latent space with \(k\) matched to the number of occupied \(F_{\rm top8}\) families for that checkpoint. The random selector permutes count-matched labels on the fitting transitions and routes later states by nearest random-label latent center. The reported quantity is the routed/global MSE ratio at \(h=100\), so values below one mean that the selected local maps improve over the same model's global rollout. Because this appendix table reports arithmetic means of per-system seed-IQM finite ratios, rather than the appendix routing table's system-level mean log effect, isolated catastrophic finite rollouts can dominate all-state cells; the deep-state slice is the cleaner partition audit.

\begin{table}[ht]
\centering
\caption{Partition-control local forecasting for the top-$8$ support-family comparison. Entries are routed/global MSE-ratio IQMs with within-system Wilcoxon/Holm counts in brackets; lower is better. The LISTA-SB row uses the matched $d_z{=}256$ control artifact.}
\label{tab:self_routing_partition_controls}
\setlength{\tabcolsep}{2.5pt}
\resizebox{\textwidth}{!}{
\begin{tabular}{@{}ll cc@{}}
\toprule
& & \multicolumn{2}{c}{$H100$} \\
\cmidrule(lr){3-4}
Model & Selector & All & Deep \\
\midrule
LISTA-SB & Oracle basin labels & $0.0176\,[15/17]$ & $0.00679\,[14/17]$ \\
 & Support family & $0.168\,[13/17]$ & $0.0374\,[14/17]$ \\
 & Latent clusters & $0.179\,[11/17]$ & $0.0471\,[13/17]$ \\
 & Random matched & $0.857\,[13/17]$ & $0.869\,[13/17]$ \\
\addlinespace
LISTA-BD & Oracle basin labels & $0.00696\,[16/17]$ & $0.00347\,[14/17]$ \\
 & Support family & $0.156\,[12/17]$ & $0.0378\,[13/17]$ \\
 & Latent clusters & $0.206\,[11/17]$ & $0.0559\,[13/17]$ \\
 & Random matched & $0.866\,[15/17]$ & $0.899\,[14/17]$ \\
\addlinespace
Sparse MLP, BD & Oracle basin labels & $0.00391\,[8/17]$ & $0.00426\,[7/17]$ \\
 & Support family & $0.222\,[7/17]$ & $0.101\,[8/17]$ \\
 & Latent clusters & $0.0139\,[13/17]$ & $0.00882\,[14/17]$ \\
 & Random matched & $0.951\,[12/17]$ & $0.978\,[10/17]$ \\
\addlinespace
Sparse MLP & Oracle basin labels & $0.00369\,[10/17]$ & $0.00366\,[10/17]$ \\
 & Support family & $0.13\,[6/17]$ & $0.0605\,[5/17]$ \\
 & Latent clusters & $0.0133\,[14/17]$ & $0.0114\,[13/17]$ \\
 & Random matched & $0.966\,[14/17]$ & $0.982\,[13/17]$ \\
\addlinespace
Dense MLP & Oracle basin labels & $10.6\,[5/17]$ & $9.32\,[4/17]$ \\
 & Support family & $1{\times}10^{3}\,[0/17]$ & $662\,[0/17]$ \\
 & Latent clusters & $1.6{\times}10^{6}\,[1/17]$ & $1.3{\times}10^{5}\,[3/17]$ \\
 & Random matched & $1.09\,[1/17]$ & $1\,[1/17]$ \\
\bottomrule
\end{tabular}}
\end{table}


The control table rules out the strongest but false interpretation of the routing result: basin labels alone are not a guaranteed route to a stable closed-loop local predictor in an arbitrary learned latent space. Oracle basin-label maps are very strong for the LISTA-style rows in this audit table, which shows that useful local centered laws exist for those learned representations. The same oracle partition is unstable for the MLP controls, including the dense no-sparsity MLP. Thus the supported conclusion is narrower and more precise: route-local centered maps can help, but the partition, the latent geometry, and the autonomous rollout dynamics must be compatible.

The non-oracle support-family rows show why \(F_{\rm top8}\) remains a meaningful mechanism variable. On basin-interior initial states, \(F_{\rm top8}\)-local routing gives large improvements for the LISTA-style rows in the table, while the dense no-sparsity MLP support-family route fails catastrophically. This supports the main claim that sparse encoders produce inspectable active-coordinate families that can act as label-free routing variables. The latent-cluster controls show that generic latent geometry can also support useful local maps, especially for sparse MLP controls, so the result should not be read as ``only support families can ever route.'' The random count-matched controls can give modest gains in some rows, but they do not reproduce the strongest deep-state support-family or latent-cluster effects and they fail for the dense MLP. The take-away is therefore comparative rather than absolute: \(F_{\rm top8}\) is a label-free, interpretable selector produced by the sparse representation, whereas oracle basins, \(k\)-means clusters, and random matched partitions are diagnostic controls rather than deployment-time support objects.

A separate routed-forecasting add-on compared the same post-hoc \(F_{\rm top8}\) family-local centered predictor against the same model's best-periodic paper-table forecasts. The fully autonomous routed version was negative for absolute forecasting superiority: it produced complete merged outputs for the retained controlled multibasin benchmark and Dysts benchmark with \(0\) failures, but for every \(d_z=256\) LISTA-family model and evaluated horizon, routed rows won \(0/15\) controlled systems and \(0/10\) Dysts systems against the same-model best-periodic baseline. A follow-up that periodically decoded and re-encoded the routed rollout every \(5\), \(10\), \(20\), or \(30\) steps also completed with \(0\) failures. Periodic refresh reduced many autonomous blow-ups and made route coverage high on the controlled systems, but it did not reverse the forecasting conclusion: every aggregate routed/best-periodic ratio remained above one, with only isolated system wins in individual cells. This negative result is important because the appendix routing table compares routed rollouts to the same model's autonomous global-\(K\) rollout, whereas the paper-table forecasting baseline is stronger because it periodically decodes and re-encodes. We therefore use support-family routing as appendix mechanism evidence that \(F_{\rm top8}\) can select useful local predictors relative to the same model's global latent rollout, not as evidence that the post-hoc routed predictor replaces best-periodic forecasting in the aggregate forecasting tables.

\section{Full per-system multibasin tables}
\label{app:fixed17_full}

Tables \ref{tab:persystem_h100}--\ref{tab:persystem_h1000} render the per-(system, candidate) effect sizes and Holm-corrected Wilcoxon \(p\)-values that underlie the \(K/N\) counts in \cref{tab:fixed17_alignment} and the forest plot in \cref{fig:fixed17_forest}. Each cell gives the per-system median paired \(\log_{10}\)-MSE difference (candidate minus matched-sampling Dense MLP, no-shrink baseline) over the common completed seeds, up to \(15\) per system, and the corresponding Holm-corrected one-sided paired Wilcoxon \(p\)-value. Bold entries indicate Holm-corrected \(p<0.05\) in favor of the candidate.

\begin{figure}[tbp]
\centering
\includegraphics[width=\linewidth]{neurips_paper_2026/fig_fixed17_forest.pdf}
\caption{Per-system effect sizes at \(H{=}1000\) versus the Dense MLP baseline. Each point is one benchmark system; the \(x\)-coordinate is the per-system median paired \(\log_{10}\)-MSE difference (candidate minus baseline) over the common completed seeds, up to \(15\) per system, with whiskers showing the \(95\%\) bootstrap interval. Filled markers clear Holm-corrected \(\alpha{=}0.05\). Numerical values are in \cref{tab:persystem_h1000}.}
\label{fig:fixed17_forest}
\end{figure}

\begin{table}[tbp]
\centering
\caption{Per-system effect sizes and Holm-corrected Wilcoxon \(p\)-values at \(H=100\) on the \(15\)-system multibasin benchmark, against the matched-sampling Dense MLP, no-shrink baseline. \(\Delta\) is the per-system median paired \(\log_{10}\)-MSE difference (candidate \( - \) baseline); \(p_{\rm H}\) is the one-sided paired Wilcoxon signed-rank \(p\)-value, Holm-corrected across retained systems within each candidate column. Bold \(\Delta\): Holm-corrected \(p<0.05\) in favor of the candidate.}
\label{tab:persystem_h100}
\setlength{\tabcolsep}{3pt}
\resizebox{\textwidth}{!}{\begin{tabular}{@{}l rr rr rr rr rr@{}}
\toprule
System & \multicolumn{2}{c}{LISTA} & \multicolumn{2}{c}{LISTA-BD} & \multicolumn{2}{c}{LISTA-SB} & \multicolumn{2}{c}{Sparse MLP, BD} & \multicolumn{2}{c}{Sparse MLP} \\
\cmidrule(lr){2-3} \cmidrule(lr){4-5} \cmidrule(lr){6-7} \cmidrule(lr){8-9} \cmidrule(lr){10-11}
 & $\Delta$ & $p_{\rm H}$ & $\Delta$ & $p_{\rm H}$ & $\Delta$ & $p_{\rm H}$ & $\Delta$ & $p_{\rm H}$ & $\Delta$ & $p_{\rm H}$ \\
\midrule
Local-linear gates & $\mathbf{-2.16}$ & $<\!10^{-3}$ & $\mathbf{-2.19}$ & $<\!10^{-3}$ & $\mathbf{-2.20}$ & $<\!10^{-3}$ & $\mathbf{-2.33}$ & $<\!10^{-3}$ & $\mathbf{-2.52}$ & $<\!10^{-3}$ \\
Transfer-gated local-linear & $\mathbf{-0.27}$ & $0.001$ & $\mathbf{-0.21}$ & $0.008$ & $\mathbf{-0.31}$ & $0.007$ & $\mathbf{-0.37}$ & $<\!10^{-3}$ & $\mathbf{-0.27}$ & $0.004$ \\
Arrested spiral & $\mathbf{-1.85}$ & $<\!10^{-3}$ & $\mathbf{-2.09}$ & $<\!10^{-3}$ & $\mathbf{-2.03}$ & $<\!10^{-3}$ & $\mathbf{-1.43}$ & $<\!10^{-3}$ & $\mathbf{-1.90}$ & $<\!10^{-3}$ \\
Asymmetric three-well & $\mathbf{-2.00}$ & $<\!10^{-3}$ & $\mathbf{-1.92}$ & $<\!10^{-3}$ & $\mathbf{-1.75}$ & $<\!10^{-3}$ & $\mathbf{-2.43}$ & $<\!10^{-3}$ & $\mathbf{-2.44}$ & $<\!10^{-3}$ \\
High-cross three-well & $\mathbf{-1.05}$ & $<\!10^{-3}$ & $\mathbf{-1.12}$ & $<\!10^{-3}$ & $\mathbf{-1.19}$ & $<\!10^{-3}$ & $\mathbf{-0.95}$ & $<\!10^{-3}$ & $\mathbf{-1.17}$ & $<\!10^{-3}$ \\
Hexagonal six-well & $\mathbf{-1.76}$ & $<\!10^{-3}$ & $\mathbf{-1.73}$ & $<\!10^{-3}$ & $\mathbf{-1.85}$ & $<\!10^{-3}$ & $\mathbf{-1.70}$ & $<\!10^{-3}$ & $\mathbf{-1.83}$ & $<\!10^{-3}$ \\
Octagonal eight-well & $\mathbf{-1.85}$ & $<\!10^{-3}$ & $\mathbf{-1.90}$ & $<\!10^{-3}$ & $\mathbf{-1.80}$ & $<\!10^{-3}$ & $\mathbf{-1.66}$ & $<\!10^{-3}$ & $\mathbf{-1.63}$ & $<\!10^{-3}$ \\
Pentagonal five-well & $\mathbf{-1.86}$ & $<\!10^{-3}$ & $\mathbf{-1.58}$ & $<\!10^{-3}$ & $\mathbf{-1.20}$ & $<\!10^{-3}$ & $\mathbf{-1.53}$ & $<\!10^{-3}$ & $\mathbf{-2.16}$ & $<\!10^{-3}$ \\
Square four-well & $\mathbf{-1.44}$ & $<\!10^{-3}$ & $\mathbf{-1.62}$ & $<\!10^{-3}$ & $\mathbf{-1.58}$ & $<\!10^{-3}$ & $\mathbf{-1.26}$ & $<\!10^{-3}$ & $\mathbf{-1.02}$ & $<\!10^{-3}$ \\
Triple-well Duffing & $\mathbf{-0.54}$ & $<\!10^{-3}$ & $\mathbf{-0.51}$ & $0.008$ & $\mathbf{-0.35}$ & $0.007$ & $\mathbf{-0.90}$ & $<\!10^{-3}$ & $\mathbf{-0.78}$ & $<\!10^{-3}$ \\
SNIC multi-attractor & $\mathbf{-2.01}$ & $<\!10^{-3}$ & $\mathbf{-2.03}$ & $<\!10^{-3}$ & $\mathbf{-1.91}$ & $<\!10^{-3}$ & $\mathbf{-2.21}$ & $<\!10^{-3}$ & $\mathbf{-2.13}$ & $<\!10^{-3}$ \\
Transition-routes four-well & $\mathbf{-2.21}$ & $<\!10^{-3}$ & $\mathbf{-1.77}$ & $<\!10^{-3}$ & $\mathbf{-1.90}$ & $<\!10^{-3}$ & $\mathbf{-2.14}$ & $<\!10^{-3}$ & $\mathbf{-2.47}$ & $<\!10^{-3}$ \\
Depth-gradient four-well & $\mathbf{-1.39}$ & $<\!10^{-3}$ & $\mathbf{-1.33}$ & $<\!10^{-3}$ & $\mathbf{-1.37}$ & $<\!10^{-3}$ & $\mathbf{-0.94}$ & $<\!10^{-3}$ & $\mathbf{-0.82}$ & $<\!10^{-3}$ \\
Diamond four-well & $\mathbf{-2.00}$ & $<\!10^{-3}$ & $\mathbf{-1.80}$ & $<\!10^{-3}$ & $\mathbf{-1.84}$ & $<\!10^{-3}$ & $\mathbf{-1.51}$ & $<\!10^{-3}$ & $\mathbf{-1.67}$ & $<\!10^{-3}$ \\
L-shaped five-well & $\mathbf{-2.23}$ & $<\!10^{-3}$ & $\mathbf{-1.94}$ & $<\!10^{-3}$ & $\mathbf{-1.90}$ & $<\!10^{-3}$ & $\mathbf{-1.85}$ & $<\!10^{-3}$ & $\mathbf{-2.23}$ & $<\!10^{-3}$ \\
\midrule
\emph{Significant/eligible} & \multicolumn{2}{c}{$K{=}15\,/\,N{=}15$} & \multicolumn{2}{c}{$K{=}15\,/\,N{=}15$} & \multicolumn{2}{c}{$K{=}15\,/\,N{=}15$} & \multicolumn{2}{c}{$K{=}15\,/\,N{=}15$} & \multicolumn{2}{c}{$K{=}15\,/\,N{=}15$} \\
\bottomrule
\end{tabular}}
\end{table}

\begin{table}[tbp]
\centering
\caption{Per-system effect sizes and Holm-corrected Wilcoxon \(p\)-values at \(H=500\) on the \(15\)-system multibasin benchmark. Format identical to \cref{tab:persystem_h100}.}
\label{tab:persystem_h500}
\setlength{\tabcolsep}{3pt}
\resizebox{\textwidth}{!}{\begin{tabular}{@{}l rr rr rr rr rr@{}}
\toprule
System & \multicolumn{2}{c}{LISTA} & \multicolumn{2}{c}{LISTA-BD} & \multicolumn{2}{c}{LISTA-SB} & \multicolumn{2}{c}{Sparse MLP, BD} & \multicolumn{2}{c}{Sparse MLP} \\
\cmidrule(lr){2-3} \cmidrule(lr){4-5} \cmidrule(lr){6-7} \cmidrule(lr){8-9} \cmidrule(lr){10-11}
 & $\Delta$ & $p_{\rm H}$ & $\Delta$ & $p_{\rm H}$ & $\Delta$ & $p_{\rm H}$ & $\Delta$ & $p_{\rm H}$ & $\Delta$ & $p_{\rm H}$ \\
\midrule
Local-linear gates & $\mathbf{-2.01}$ & $<\!10^{-3}$ & $\mathbf{-2.24}$ & $<\!10^{-3}$ & $\mathbf{-2.27}$ & $<\!10^{-3}$ & $\mathbf{-2.29}$ & $<\!10^{-3}$ & $\mathbf{-2.52}$ & $<\!10^{-3}$ \\
Transfer-gated local-linear & $\mathbf{-0.08}$ & $0.013$ & $\mathbf{+1.19}$ & $0.008$ & $\mathbf{-0.21}$ & $<\!10^{-3}$ & $\mathbf{-0.17}$ & $<\!10^{-3}$ & $\mathbf{-0.18}$ & $<\!10^{-3}$ \\
Arrested spiral & $\mathbf{-2.26}$ & $<\!10^{-3}$ & $\mathbf{-2.57}$ & $<\!10^{-3}$ & $\mathbf{-2.57}$ & $<\!10^{-3}$ & $\mathbf{-1.72}$ & $<\!10^{-3}$ & $\mathbf{-2.33}$ & $<\!10^{-3}$ \\
Asymmetric three-well & $\mathbf{-2.27}$ & $<\!10^{-3}$ & $\mathbf{-2.12}$ & $<\!10^{-3}$ & $\mathbf{-1.94}$ & $<\!10^{-3}$ & $\mathbf{-3.11}$ & $<\!10^{-3}$ & $\mathbf{-3.15}$ & $<\!10^{-3}$ \\
High-cross three-well & $\mathbf{-0.93}$ & $<\!10^{-3}$ & $\mathbf{-0.97}$ & $<\!10^{-3}$ & $\mathbf{-1.03}$ & $<\!10^{-3}$ & $\mathbf{-0.86}$ & $<\!10^{-3}$ & $\mathbf{-1.26}$ & $<\!10^{-3}$ \\
Hexagonal six-well & $\mathbf{-2.01}$ & $<\!10^{-3}$ & $\mathbf{-2.01}$ & $<\!10^{-3}$ & $\mathbf{-2.18}$ & $<\!10^{-3}$ & $\mathbf{-1.95}$ & $<\!10^{-3}$ & $\mathbf{-2.16}$ & $<\!10^{-3}$ \\
Octagonal eight-well & $\mathbf{-1.94}$ & $<\!10^{-3}$ & $\mathbf{-2.09}$ & $<\!10^{-3}$ & $\mathbf{-1.91}$ & $<\!10^{-3}$ & $\mathbf{-1.78}$ & $<\!10^{-3}$ & $\mathbf{-1.74}$ & $<\!10^{-3}$ \\
Pentagonal five-well & $\mathbf{-2.48}$ & $<\!10^{-3}$ & $\mathbf{-2.10}$ & $<\!10^{-3}$ & $\mathbf{-1.61}$ & $<\!10^{-3}$ & $\mathbf{-2.02}$ & $<\!10^{-3}$ & $\mathbf{-2.77}$ & $<\!10^{-3}$ \\
Square four-well & $\mathbf{-1.60}$ & $<\!10^{-3}$ & $\mathbf{-1.89}$ & $<\!10^{-3}$ & $\mathbf{-1.81}$ & $<\!10^{-3}$ & $\mathbf{-1.55}$ & $<\!10^{-3}$ & $\mathbf{-1.17}$ & $<\!10^{-3}$ \\
Triple-well Duffing & $\mathbf{-0.83}$ & $<\!10^{-3}$ & $\mathbf{-0.81}$ & $<\!10^{-3}$ & $\mathbf{-0.65}$ & $<\!10^{-3}$ & $\mathbf{-1.10}$ & $<\!10^{-3}$ & $\mathbf{-1.05}$ & $<\!10^{-3}$ \\
SNIC multi-attractor & $\mathbf{-2.30}$ & $<\!10^{-3}$ & $\mathbf{-2.45}$ & $<\!10^{-3}$ & $\mathbf{-2.31}$ & $<\!10^{-3}$ & $\mathbf{-2.56}$ & $<\!10^{-3}$ & $\mathbf{-2.49}$ & $<\!10^{-3}$ \\
Transition-routes four-well & $\mathbf{-1.73}$ & $<\!10^{-3}$ & $\mathbf{-1.52}$ & $<\!10^{-3}$ & $\mathbf{-1.49}$ & $<\!10^{-3}$ & $\mathbf{-1.55}$ & $<\!10^{-3}$ & $\mathbf{-1.86}$ & $<\!10^{-3}$ \\
Depth-gradient four-well & $\mathbf{-1.63}$ & $<\!10^{-3}$ & $\mathbf{-1.45}$ & $<\!10^{-3}$ & $\mathbf{-1.47}$ & $<\!10^{-3}$ & $\mathbf{-1.17}$ & $<\!10^{-3}$ & $\mathbf{-0.92}$ & $<\!10^{-3}$ \\
Diamond four-well & $\mathbf{-1.77}$ & $<\!10^{-3}$ & $\mathbf{-1.58}$ & $<\!10^{-3}$ & $\mathbf{-1.40}$ & $<\!10^{-3}$ & $\mathbf{-1.28}$ & $<\!10^{-3}$ & $\mathbf{-1.75}$ & $<\!10^{-3}$ \\
L-shaped five-well & $\mathbf{-2.98}$ & $<\!10^{-3}$ & $\mathbf{-2.04}$ & $<\!10^{-3}$ & $\mathbf{-2.03}$ & $<\!10^{-3}$ & $\mathbf{-2.27}$ & $<\!10^{-3}$ & $\mathbf{-2.96}$ & $<\!10^{-3}$ \\
\midrule
\emph{Significant/eligible} & \multicolumn{2}{c}{$K{=}15\,/\,N{=}15$} & \multicolumn{2}{c}{$K{=}15\,/\,N{=}15$} & \multicolumn{2}{c}{$K{=}15\,/\,N{=}15$} & \multicolumn{2}{c}{$K{=}15\,/\,N{=}15$} & \multicolumn{2}{c}{$K{=}15\,/\,N{=}15$} \\
\bottomrule
\end{tabular}}
\end{table}

\begin{table}[tbp]
\centering
\caption{Per-system effect sizes and Holm-corrected Wilcoxon \(p\)-values at \(H=1000\) on the \(15\)-system multibasin benchmark. This is the numerical version of the forest plot in \cref{fig:fixed17_forest}. Format identical to \cref{tab:persystem_h100}.}
\label{tab:persystem_h1000}
\setlength{\tabcolsep}{3pt}
\resizebox{\textwidth}{!}{\begin{tabular}{@{}l rr rr rr rr@{}}
\toprule
System & \multicolumn{2}{c}{LISTA-BD} & \multicolumn{2}{c}{LISTA-SB} & \multicolumn{2}{c}{Sparse MLP, BD} & \multicolumn{2}{c}{Sparse MLP} \\
\cmidrule(lr){2-3} \cmidrule(lr){4-5} \cmidrule(lr){6-7} \cmidrule(lr){8-9}
 & $\Delta$ & $p_{\rm H}$ & $\Delta$ & $p_{\rm H}$ & $\Delta$ & $p_{\rm H}$ & $\Delta$ & $p_{\rm H}$ \\
\midrule
checkerboard\_potential & $-2.78$ & $0.216$ & $\mathbf{-2.55}$ & $0.017$ & $\mathbf{-2.90}$ & $0.017$ & $\mathbf{-2.61}$ & $0.017$ \\
snic\_multi & $\mathbf{-2.37}$ & $0.017$ & $\mathbf{-1.88}$ & $0.017$ & $\mathbf{-2.42}$ & $0.017$ & $\mathbf{-2.39}$ & $0.017$ \\
gated\_local\_linear & $\mathbf{-2.23}$ & $0.017$ & $\mathbf{-2.21}$ & $0.017$ & $\mathbf{-2.38}$ & $0.017$ & $\mathbf{-2.70}$ & $0.017$ \\
var\_l\_shape\_5 & $\mathbf{-1.89}$ & $0.017$ & $\mathbf{-2.07}$ & $0.017$ & $\mathbf{-2.45}$ & $0.017$ & $\mathbf{-2.55}$ & $0.017$ \\
cal\_hexagon\_6 & $\mathbf{-1.92}$ & $0.017$ & $\mathbf{-2.01}$ & $0.017$ & $\mathbf{-2.14}$ & $0.017$ & $\mathbf{-2.14}$ & $0.017$ \\
arrested\_spiral & $\mathbf{-2.14}$ & $0.017$ & $-1.81$ & $0.129$ & $\mathbf{-2.08}$ & $0.017$ & $\mathbf{-2.05}$ & $0.017$ \\
cal\_square\_4 & $\mathbf{-2.06}$ & $0.017$ & $\mathbf{-1.28}$ & $0.017$ & $\mathbf{-2.26}$ & $0.017$ & $\mathbf{-2.04}$ & $0.017$ \\
cal\_pentagon\_5 & $\mathbf{-1.70}$ & $0.017$ & $\mathbf{-1.31}$ & $0.017$ & $\mathbf{-1.98}$ & $0.017$ & $\mathbf{-2.83}$ & $0.017$ \\
cal\_octagon\_8 & $\mathbf{-2.02}$ & $0.017$ & $\mathbf{-1.79}$ & $0.017$ & $\mathbf{-1.87}$ & $0.017$ & $\mathbf{-1.65}$ & $0.017$ \\
cal\_asymmetric\_3 & $\mathbf{-1.75}$ & $0.017$ & $\mathbf{-1.48}$ & $0.017$ & $\mathbf{-1.77}$ & $0.017$ & $\mathbf{-3.60}$ & $0.017$ \\
multiwell\_strong\_transition & $-1.63$ & $0.084$ & $-2.42$ & $0.500$ & $\mathbf{-1.70}$ & $0.042$ & $\mathbf{-1.75}$ & $0.042$ \\
transition\_routes\_4 & $\mathbf{-1.35}$ & $0.017$ & $\mathbf{-1.72}$ & $0.017$ & $\mathbf{-1.55}$ & $0.017$ & $\mathbf{-1.63}$ & $0.017$ \\
var\_diamond\_4 & $\mathbf{-1.57}$ & $0.017$ & $\mathbf{-1.32}$ & $0.017$ & $\mathbf{-1.57}$ & $0.017$ & $\mathbf{-2.00}$ & $0.017$ \\
var\_depth\_gradient\_4 & $\mathbf{-2.00}$ & $0.017$ & $\mathbf{-1.05}$ & $0.017$ & $\mathbf{-2.27}$ & $0.017$ & $\mathbf{-0.55}$ & $0.017$ \\
cal\_high\_cross\_3 & $\mathbf{-1.19}$ & $0.017$ & $\mathbf{-1.24}$ & $0.017$ & $\mathbf{-1.07}$ & $0.017$ & $\mathbf{-1.38}$ & $0.017$ \\
duffing\_triple\_well & $\mathbf{-0.81}$ & $0.017$ & $\mathbf{-0.75}$ & $0.017$ & $\mathbf{-1.21}$ & $0.017$ & $\mathbf{-1.09}$ & $0.017$ \\
gated\_transfer\_linear & $\mathbf{-0.15}$ & $0.017$ & $\mathbf{-0.13}$ & $0.017$ & $\mathbf{-0.16}$ & $0.017$ & $\mathbf{-0.17}$ & $0.017$ \\
\midrule
\emph{$K/17$ Holm $p<0.05$} & \multicolumn{2}{c}{$K{=}15\,/\,N{=}17$} & \multicolumn{2}{c}{$K{=}15\,/\,N{=}17$} & \multicolumn{2}{c}{$K{=}17\,/\,N{=}17$} & \multicolumn{2}{c}{$K{=}17\,/\,N{=}17$} \\
\bottomrule
\end{tabular}}
\end{table}

\subsection*{Per-system audit of the interpretability columns}

\paragraph{Why the wrong-support denominator is \(11\) and not \(15\).}
The wrong-support ablation builds, for each system, a dictionary mapping each basin to its canonical support mask, defined as the most-common exact support among the deep-slice states of that basin. The ablated rollout draws a mask from a basin different from the state's own basin and holds that wrong mask fixed, so the test requires at least two distinct deep-slice canonical masks. Four of the \(15\) retained benchmark systems, the arrested spiral, asymmetric three-well, triple-well Duffing, and L-shaped five-well systems, have multiple basins by design but a global top-quartile deep slice concentrated in a single basin. The wrong-support ablation is undefined on those four systems by construction, and the corresponding \(K/N\) denominator in \cref{tab:fixed17_alignment} and \cref{tab:persystem_freeze_h1,tab:persystem_freeze_h20} is \(11\) instead of \(15\). A per-basin top-quartile depth re-evaluation is the appendix robustness analysis for the alternative basin-stratified coverage question, but the global-slice numbers remain the source of the current main table.

Tables \ref{tab:persystem_size}--\ref{tab:persystem_freeze_h20} render the full per-(system, candidate) breakdown of the retained active-count diagnostic and the three directional interpretability columns of \cref{tab:fixed17_alignment}. The alternative direction is ``candidate \(<\) baseline'' for \(|S_{\rm abs}|\) and \(H(B\mid F_{\rm abs})\), and ``candidate \(>\) baseline'' for the wrong-support ablation ratios. The main-text family-count column is intentionally not included in these one-sided tests because \(|F_{\rm abs}|\) is a non-directional count.

\begin{table}[tbp]
\centering
\caption{Per-system effect sizes and Holm-corrected Wilcoxon \(p\)-values for \(|S_{\rm abs}|\), the mean active-coordinate count under the absolute-threshold support on the deep slice. Alternative: candidate \(|S_{\rm abs}|<\) baseline \(|S_{\rm abs}|\).}
\label{tab:persystem_size}
\setlength{\tabcolsep}{3pt}
\resizebox{\textwidth}{!}{\begin{tabular}{@{}l rr rr rr rr@{}}
\toprule
System & \multicolumn{2}{c}{LISTA-BD} & \multicolumn{2}{c}{LISTA-SB} & \multicolumn{2}{c}{Sparse MLP, BD} & \multicolumn{2}{c}{Sparse MLP} \\
\cmidrule(lr){2-3} \cmidrule(lr){4-5} \cmidrule(lr){6-7} \cmidrule(lr){8-9}
 & $\Delta$ & $p_{\rm H}$ & $\Delta$ & $p_{\rm H}$ & $\Delta$ & $p_{\rm H}$ & $\Delta$ & $p_{\rm H}$ \\
\midrule
snic\_multi & $\mathbf{-174.91}$ & $0.017$ & $\mathbf{-233.11}$ & $0.017$ & $\mathbf{-171.26}$ & $0.017$ & $\mathbf{-171.16}$ & $0.017$ \\
gated\_local\_linear & $\mathbf{-167.78}$ & $0.017$ & $\mathbf{-229.59}$ & $0.017$ & $\mathbf{-160.43}$ & $0.017$ & $\mathbf{-160.28}$ & $0.017$ \\
cal\_asymmetric\_3 & $\mathbf{-158.00}$ & $0.017$ & $\mathbf{-229.02}$ & $0.017$ & $\mathbf{-159.00}$ & $0.017$ & $\mathbf{-162.00}$ & $0.017$ \\
gated\_transfer\_linear & $\mathbf{-154.26}$ & $0.017$ & $\mathbf{-227.24}$ & $0.017$ & $\mathbf{-157.98}$ & $0.017$ & $\mathbf{-158.41}$ & $0.017$ \\
transition\_routes\_4 & $\mathbf{-162.97}$ & $0.017$ & $\mathbf{-228.84}$ & $0.017$ & $\mathbf{-152.28}$ & $0.017$ & $\mathbf{-152.68}$ & $0.017$ \\
duffing\_triple\_well & $\mathbf{-163.37}$ & $0.017$ & $\mathbf{-230.04}$ & $0.017$ & $\mathbf{-149.16}$ & $0.017$ & $\mathbf{-149.16}$ & $0.017$ \\
var\_diamond\_4 & $\mathbf{-147.30}$ & $0.017$ & $\mathbf{-226.48}$ & $0.017$ & $\mathbf{-155.50}$ & $0.017$ & $\mathbf{-155.61}$ & $0.017$ \\
arrested\_spiral & $\mathbf{-149.91}$ & $0.017$ & $\mathbf{-225.92}$ & $0.017$ & $\mathbf{-154.04}$ & $0.017$ & $\mathbf{-154.19}$ & $0.017$ \\
cal\_pentagon\_5 & $\mathbf{-156.14}$ & $0.017$ & $\mathbf{-225.99}$ & $0.017$ & $\mathbf{-142.92}$ & $0.017$ & $\mathbf{-142.86}$ & $0.017$ \\
multiwell\_strong\_transition & $\mathbf{-137.88}$ & $0.017$ & $\mathbf{-225.31}$ & $0.017$ & $\mathbf{-149.39}$ & $0.017$ & $\mathbf{-149.34}$ & $0.017$ \\
cal\_hexagon\_6 & $\mathbf{-153.12}$ & $0.017$ & $\mathbf{-227.34}$ & $0.017$ & $\mathbf{-143.44}$ & $0.017$ & $\mathbf{-143.44}$ & $0.017$ \\
var\_l\_shape\_5 & $\mathbf{-151.37}$ & $0.017$ & $\mathbf{-227.92}$ & $0.017$ & $\mathbf{-143.60}$ & $0.017$ & $\mathbf{-143.42}$ & $0.017$ \\
cal\_square\_4 & $\mathbf{-152.42}$ & $0.017$ & $\mathbf{-225.23}$ & $0.017$ & $\mathbf{-141.23}$ & $0.017$ & $\mathbf{-141.23}$ & $0.017$ \\
var\_depth\_gradient\_4 & $\mathbf{-150.24}$ & $0.017$ & $\mathbf{-225.29}$ & $0.017$ & $\mathbf{-140.58}$ & $0.017$ & $\mathbf{-140.24}$ & $0.017$ \\
cal\_octagon\_8 & $\mathbf{-146.49}$ & $0.017$ & $\mathbf{-226.48}$ & $0.017$ & $\mathbf{-143.17}$ & $0.017$ & $\mathbf{-143.17}$ & $0.017$ \\
checkerboard\_potential & $\mathbf{-143.32}$ & $0.017$ & $\mathbf{-221.69}$ & $0.017$ & $\mathbf{-138.58}$ & $0.017$ & $\mathbf{-138.58}$ & $0.017$ \\
cal\_high\_cross\_3 & $\mathbf{-144.04}$ & $0.017$ & $\mathbf{-226.00}$ & $0.017$ & $\mathbf{-134.51}$ & $0.017$ & $\mathbf{-135.50}$ & $0.017$ \\
\midrule
\emph{$K/17$ Holm $p<0.05$} & \multicolumn{2}{c}{$K{=}17\,/\,N{=}17$} & \multicolumn{2}{c}{$K{=}17\,/\,N{=}17$} & \multicolumn{2}{c}{$K{=}17\,/\,N{=}17$} & \multicolumn{2}{c}{$K{=}17\,/\,N{=}17$} \\
\bottomrule
\end{tabular}}
\end{table}

\begin{table}[tbp]
\centering
\caption{Per-system effect sizes and Holm-corrected Wilcoxon \(p\)-values for \(H(B\mid F_{\rm abs})\) on the deep slice. Alternative: candidate \(H(B\mid F_{\rm abs})<\) baseline.}
\label{tab:persystem_hbf}
\setlength{\tabcolsep}{3pt}
\resizebox{\textwidth}{!}{\begin{tabular}{@{}l rr rr rr rr rr@{}}
\toprule
System & \multicolumn{2}{c}{LISTA} & \multicolumn{2}{c}{LISTA-BD} & \multicolumn{2}{c}{LISTA-SB} & \multicolumn{2}{c}{Sparse MLP, BD} & \multicolumn{2}{c}{Sparse MLP} \\
\cmidrule(lr){2-3} \cmidrule(lr){4-5} \cmidrule(lr){6-7} \cmidrule(lr){8-9} \cmidrule(lr){10-11}
 & $\Delta$ & $p_{\rm H}$ & $\Delta$ & $p_{\rm H}$ & $\Delta$ & $p_{\rm H}$ & $\Delta$ & $p_{\rm H}$ & $\Delta$ & $p_{\rm H}$ \\
\midrule
Local-linear gates & $\mathbf{-1.09}$ & $<\!10^{-3}$ & $\mathbf{-1.09}$ & $<\!10^{-3}$ & $\mathbf{-1.09}$ & $<\!10^{-3}$ & $\mathbf{-1.09}$ & $<\!10^{-3}$ & $\mathbf{-1.09}$ & $<\!10^{-3}$ \\
Transfer-gated local-linear & $\mathbf{-1.10}$ & $<\!10^{-3}$ & $\mathbf{-1.10}$ & $<\!10^{-3}$ & $\mathbf{-1.10}$ & $<\!10^{-3}$ & $\mathbf{-1.10}$ & $<\!10^{-3}$ & $\mathbf{-1.10}$ & $<\!10^{-3}$ \\
Arrested spiral & $\mathbf{-1.04}$ & $0.001$ & $\mathbf{-0.92}$ & $0.001$ & $\mathbf{-0.98}$ & $0.001$ & $\mathbf{-0.27}$ & $0.026$ & $\mathbf{-0.98}$ & $0.001$ \\
Asymmetric three-well & $\mathbf{-1.09}$ & $<\!10^{-3}$ & $\mathbf{-1.09}$ & $<\!10^{-3}$ & $\mathbf{-1.09}$ & $<\!10^{-3}$ & $\mathbf{-1.09}$ & $<\!10^{-3}$ & $\mathbf{-1.09}$ & $<\!10^{-3}$ \\
High-cross three-well & $\mathbf{-0.95}$ & $<\!10^{-3}$ & $\mathbf{-0.79}$ & $0.001$ & $\mathbf{-1.01}$ & $<\!10^{-3}$ & $\mathbf{-0.20}$ & $0.026$ & $\mathbf{-0.70}$ & $0.001$ \\
Hexagonal six-well & $\mathbf{-1.40}$ & $0.001$ & $\mathbf{-1.31}$ & $0.001$ & $\mathbf{-1.17}$ & $0.001$ & $\mathbf{-0.77}$ & $0.001$ & $\mathbf{-0.82}$ & $0.001$ \\
Octagonal eight-well & $\mathbf{-1.34}$ & $0.001$ & $\mathbf{-1.25}$ & $0.001$ & $\mathbf{-1.34}$ & $0.001$ & $\mathbf{-1.08}$ & $0.001$ & $\mathbf{-1.09}$ & $0.001$ \\
Pentagonal five-well & $\mathbf{-1.57}$ & $<\!10^{-3}$ & $\mathbf{-1.51}$ & $<\!10^{-3}$ & $\mathbf{-1.41}$ & $0.001$ & $\mathbf{-0.95}$ & $0.001$ & $\mathbf{-1.34}$ & $0.001$ \\
Square four-well & $\mathbf{-1.34}$ & $<\!10^{-3}$ & $\mathbf{-1.36}$ & $<\!10^{-3}$ & $\mathbf{-1.34}$ & $<\!10^{-3}$ & $\mathbf{-1.32}$ & $<\!10^{-3}$ & $\mathbf{-1.36}$ & $<\!10^{-3}$ \\
Triple-well Duffing & $+0.00$ & -- & $+0.00$ & -- & $+0.00$ & -- & $+0.00$ & -- & $+0.00$ & -- \\
SNIC multi-attractor & $\mathbf{-1.09}$ & $<\!10^{-3}$ & $\mathbf{-1.09}$ & $<\!10^{-3}$ & $\mathbf{-1.09}$ & $<\!10^{-3}$ & $\mathbf{-1.09}$ & $<\!10^{-3}$ & $\mathbf{-1.09}$ & $<\!10^{-3}$ \\
Transition-routes four-well & $\mathbf{-1.39}$ & $<\!10^{-3}$ & $\mathbf{-1.39}$ & $<\!10^{-3}$ & $\mathbf{-1.39}$ & $<\!10^{-3}$ & $\mathbf{-1.39}$ & $<\!10^{-3}$ & $\mathbf{-1.39}$ & $<\!10^{-3}$ \\
Depth-gradient four-well & $\mathbf{-1.30}$ & $<\!10^{-3}$ & $\mathbf{-1.34}$ & $<\!10^{-3}$ & $\mathbf{-1.36}$ & $<\!10^{-3}$ & $\mathbf{-1.24}$ & $0.001$ & $\mathbf{-1.36}$ & $<\!10^{-3}$ \\
Diamond four-well & $\mathbf{-1.33}$ & $<\!10^{-3}$ & $\mathbf{-1.33}$ & $<\!10^{-3}$ & $\mathbf{-1.36}$ & $<\!10^{-3}$ & $\mathbf{-1.36}$ & $<\!10^{-3}$ & $\mathbf{-1.36}$ & $<\!10^{-3}$ \\
L-shaped five-well & $\mathbf{-1.01}$ & $0.001$ & $\mathbf{-0.96}$ & $0.001$ & $\mathbf{-0.99}$ & $0.001$ & $\mathbf{-0.92}$ & $<\!10^{-3}$ & $\mathbf{-0.92}$ & $<\!10^{-3}$ \\
\midrule
\emph{Significant/eligible} & \multicolumn{2}{c}{$K{=}14\,/\,N{=}14$} & \multicolumn{2}{c}{$K{=}14\,/\,N{=}14$} & \multicolumn{2}{c}{$K{=}14\,/\,N{=}14$} & \multicolumn{2}{c}{$K{=}14\,/\,N{=}14$} & \multicolumn{2}{c}{$K{=}14\,/\,N{=}14$} \\
\bottomrule
\end{tabular}}
\end{table}

\begin{table}[tbp]
\centering
\caption{Per-system effect sizes and Holm-corrected Wilcoxon \(p\)-values for the wrong-support-over-base MSE ratio at \(H1\) on the deep slice. Alternative: candidate ratio \(>\) baseline.}
\label{tab:persystem_freeze_h1}
\setlength{\tabcolsep}{3pt}
\resizebox{\textwidth}{!}{\begin{tabular}{@{}l rr rr rr rr rr@{}}
\toprule
System & \multicolumn{2}{c}{LISTA} & \multicolumn{2}{c}{LISTA-BD} & \multicolumn{2}{c}{LISTA-SB} & \multicolumn{2}{c}{Sparse MLP, BD} & \multicolumn{2}{c}{Sparse MLP} \\
\cmidrule(lr){2-3} \cmidrule(lr){4-5} \cmidrule(lr){6-7} \cmidrule(lr){8-9} \cmidrule(lr){10-11}
 & $\Delta$ & $p_{\rm H}$ & $\Delta$ & $p_{\rm H}$ & $\Delta$ & $p_{\rm H}$ & $\Delta$ & $p_{\rm H}$ & $\Delta$ & $p_{\rm H}$ \\
\midrule
multiwell\_strong\_transition & $\mathbf{+761.22}$ & $<\!10^{-3}$ & $\mathbf{+506.16}$ & $<\!10^{-3}$ & $\mathbf{+463.88}$ & $<\!10^{-3}$ & $\mathbf{+1346.70}$ & $<\!10^{-3}$ & $\mathbf{+1346.70}$ & $<\!10^{-3}$ \\
checkerboard\_potential & $\mathbf{+915.51}$ & $<\!10^{-3}$ & $\mathbf{+1165.01}$ & $<\!10^{-3}$ & $\mathbf{+1038.03}$ & $<\!10^{-3}$ & $\mathbf{+2847.10}$ & $<\!10^{-3}$ & $\mathbf{+2611.92}$ & $<\!10^{-3}$ \\
cal\_high\_cross\_3 & $\mathbf{+4977.62}$ & $<\!10^{-3}$ & $\mathbf{+1471.93}$ & $<\!10^{-3}$ & $\mathbf{+6293.08}$ & $<\!10^{-3}$ & $\mathbf{+2932.96}$ & $<\!10^{-3}$ & $\mathbf{+3155.16}$ & $<\!10^{-3}$ \\
cal\_pentagon\_5 & $\mathbf{+5581.15}$ & $<\!10^{-3}$ & $\mathbf{+5222.11}$ & $<\!10^{-3}$ & $\mathbf{+4328.40}$ & $<\!10^{-3}$ & $\mathbf{+1359.93}$ & $<\!10^{-3}$ & $\mathbf{+1231.67}$ & $<\!10^{-3}$ \\
var\_depth\_gradient\_4 & $\mathbf{+4678.47}$ & $<\!10^{-3}$ & $\mathbf{+4566.67}$ & $<\!10^{-3}$ & $\mathbf{+8591.21}$ & $<\!10^{-3}$ & $\mathbf{+1036.02}$ & $<\!10^{-3}$ & $\mathbf{+1036.02}$ & $<\!10^{-3}$ \\
cal\_octagon\_8 & $\mathbf{+5016.69}$ & $<\!10^{-3}$ & $\mathbf{+6473.72}$ & $<\!10^{-3}$ & $\mathbf{+8701.56}$ & $<\!10^{-3}$ & $\mathbf{+3832.77}$ & $<\!10^{-3}$ & $\mathbf{+3782.78}$ & $<\!10^{-3}$ \\
gated\_transfer\_linear & $\mathbf{+3670.47}$ & $<\!10^{-3}$ & $\mathbf{+4362.02}$ & $<\!10^{-3}$ & $\mathbf{+5438.15}$ & $<\!10^{-3}$ & $\mathbf{+9156.48}$ & $<\!10^{-3}$ & $\mathbf{+7648.11}$ & $<\!10^{-3}$ \\
cal\_square\_4 & $\mathbf{+10647.42}$ & $<\!10^{-3}$ & $\mathbf{+13062.48}$ & $<\!10^{-3}$ & $\mathbf{+17067.95}$ & $<\!10^{-3}$ & $\mathbf{+5820.10}$ & $<\!10^{-3}$ & $\mathbf{+5395.17}$ & $<\!10^{-3}$ \\
var\_diamond\_4 & $\mathbf{+14742.42}$ & $<\!10^{-3}$ & $\mathbf{+5461.40}$ & $<\!10^{-3}$ & $\mathbf{+16175.88}$ & $<\!10^{-3}$ & $\mathbf{+8872.06}$ & $<\!10^{-3}$ & $\mathbf{+10718.65}$ & $<\!10^{-3}$ \\
cal\_hexagon\_6 & $\mathbf{+13304.86}$ & $<\!10^{-3}$ & $\mathbf{+7892.75}$ & $<\!10^{-3}$ & $\mathbf{+9778.76}$ & $<\!10^{-3}$ & $\mathbf{+12232.14}$ & $<\!10^{-3}$ & $\mathbf{+12232.14}$ & $<\!10^{-3}$ \\
snic\_multi & $\mathbf{+27930.66}$ & $<\!10^{-3}$ & $\mathbf{+18483.93}$ & $<\!10^{-3}$ & $\mathbf{+23485.14}$ & $<\!10^{-3}$ & $\mathbf{+209357.97}$ & $<\!10^{-3}$ & $\mathbf{+322696.68}$ & $<\!10^{-3}$ \\
transition\_routes\_4 & $\mathbf{+70692.42}$ & $<\!10^{-3}$ & $\mathbf{+110074.20}$ & $<\!10^{-3}$ & $\mathbf{+76351.70}$ & $<\!10^{-3}$ & $\mathbf{+32116.22}$ & $<\!10^{-3}$ & $\mathbf{+27463.47}$ & $<\!10^{-3}$ \\
gated\_local\_linear & $\mathbf{+77655.29}$ & $<\!10^{-3}$ & $\mathbf{+55127.58}$ & $<\!10^{-3}$ & $\mathbf{+112929.91}$ & $<\!10^{-3}$ & $\mathbf{+281342.90}$ & $<\!10^{-3}$ & $\mathbf{+362761.28}$ & $<\!10^{-3}$ \\
\midrule
\emph{$K/17$ Holm $p<0.05$} & \multicolumn{2}{c}{$K{=}13\,/\,N{=}13$} & \multicolumn{2}{c}{$K{=}13\,/\,N{=}13$} & \multicolumn{2}{c}{$K{=}13\,/\,N{=}13$} & \multicolumn{2}{c}{$K{=}13\,/\,N{=}13$} & \multicolumn{2}{c}{$K{=}13\,/\,N{=}13$} \\
\bottomrule
\end{tabular}}
\end{table}

\begin{table}[tbp]
\centering
\caption{Per-system effect sizes and Holm-corrected Wilcoxon \(p\)-values for the wrong-support-over-base MSE ratio at \(H20\). Alternative: candidate \(>\) baseline.}
\label{tab:persystem_freeze_h20}
\setlength{\tabcolsep}{3pt}
\resizebox{\textwidth}{!}{\begin{tabular}{@{}l rr rr rr rr rr@{}}
\toprule
System & \multicolumn{2}{c}{LISTA} & \multicolumn{2}{c}{LISTA-BD} & \multicolumn{2}{c}{LISTA-SB} & \multicolumn{2}{c}{Sparse MLP, BD} & \multicolumn{2}{c}{Sparse MLP} \\
\cmidrule(lr){2-3} \cmidrule(lr){4-5} \cmidrule(lr){6-7} \cmidrule(lr){8-9} \cmidrule(lr){10-11}
 & $\Delta$ & $p_{\rm H}$ & $\Delta$ & $p_{\rm H}$ & $\Delta$ & $p_{\rm H}$ & $\Delta$ & $p_{\rm H}$ & $\Delta$ & $p_{\rm H}$ \\
\midrule
cal\_high\_cross\_3 & $+1.91$ & $0.076$ & $-0.338$ & $0.756$ & $+0.790$ & $0.094$ & $\mathbf{+3.46}$ & $0.036$ & $\mathbf{+13.9}$ & $0.008$ \\
var\_depth\_gradient\_4 & $\mathbf{+4.38}$ & $0.035$ & $\mathbf{+0.919}$ & $0.018$ & $\mathbf{+3.01}$ & $<\!10^{-3}$ & $\mathbf{+13.6}$ & $<\!10^{-3}$ & $\mathbf{+13.6}$ & $<\!10^{-3}$ \\
cal\_octagon\_8 & $\mathbf{+9.78}$ & $<\!10^{-3}$ & $\mathbf{+4.45}$ & $<\!10^{-3}$ & $\mathbf{+2.99}$ & $<\!10^{-3}$ & $\mathbf{+6.57}$ & $<\!10^{-3}$ & $\mathbf{+8.59}$ & $<\!10^{-3}$ \\
cal\_square\_4 & $\mathbf{+7.85}$ & $<\!10^{-3}$ & $\mathbf{+9.39}$ & $<\!10^{-3}$ & $\mathbf{+3.43}$ & $<\!10^{-3}$ & $\mathbf{+10.6}$ & $<\!10^{-3}$ & $\mathbf{+10.6}$ & $<\!10^{-3}$ \\
cal\_pentagon\_5 & $\mathbf{+11.0}$ & $<\!10^{-3}$ & $\mathbf{+12.3}$ & $<\!10^{-3}$ & $\mathbf{+7.70}$ & $0.001$ & $\mathbf{+11.0}$ & $<\!10^{-3}$ & $\mathbf{+8.97}$ & $<\!10^{-3}$ \\
gated\_transfer\_linear & $\mathbf{+8.94}$ & $0.002$ & $\mathbf{+22.3}$ & $0.001$ & $\mathbf{+14.5}$ & $<\!10^{-3}$ & $\mathbf{+27.5}$ & $<\!10^{-3}$ & $\mathbf{+22.1}$ & $<\!10^{-3}$ \\
cal\_hexagon\_6 & $\mathbf{+29.3}$ & $<\!10^{-3}$ & $\mathbf{+26.0}$ & $<\!10^{-3}$ & $\mathbf{+7.23}$ & $<\!10^{-3}$ & $\mathbf{+45.8}$ & $<\!10^{-3}$ & $\mathbf{+45.8}$ & $<\!10^{-3}$ \\
var\_diamond\_4 & $\mathbf{+55.7}$ & $<\!10^{-3}$ & $\mathbf{+1.75}$ & $0.004$ & $\mathbf{+40.9}$ & $<\!10^{-3}$ & $\mathbf{+49.4}$ & $<\!10^{-3}$ & $\mathbf{+49.2}$ & $<\!10^{-3}$ \\
transition\_routes\_4 & $\mathbf{+308}$ & $<\!10^{-3}$ & $\mathbf{+176}$ & $<\!10^{-3}$ & $\mathbf{+214}$ & $<\!10^{-3}$ & $\mathbf{+209}$ & $<\!10^{-3}$ & $\mathbf{+154}$ & $<\!10^{-3}$ \\
snic\_multi & $\mathbf{+214}$ & $<\!10^{-3}$ & $\mathbf{+282}$ & $<\!10^{-3}$ & $\mathbf{+186}$ & $<\!10^{-3}$ & $\mathbf{+573}$ & $<\!10^{-3}$ & $\mathbf{+626}$ & $<\!10^{-3}$ \\
gated\_local\_linear & $\mathbf{+568}$ & $<\!10^{-3}$ & $\mathbf{+362}$ & $<\!10^{-3}$ & $\mathbf{+1.17{\times}10^{3}}$ & $<\!10^{-3}$ & $\mathbf{+2.95{\times}10^{3}}$ & $<\!10^{-3}$ & $\mathbf{+4.07{\times}10^{3}}$ & $<\!10^{-3}$ \\
\midrule
\emph{$K/15$ Holm $p<0.05$} & \multicolumn{2}{c}{$K{=}10\,/\,N{=}11$} & \multicolumn{2}{c}{$K{=}10\,/\,N{=}11$} & \multicolumn{2}{c}{$K{=}10\,/\,N{=}11$} & \multicolumn{2}{c}{$K{=}11\,/\,N{=}11$} & \multicolumn{2}{c}{$K{=}11\,/\,N{=}11$} \\
\bottomrule
\end{tabular}}
\end{table}

\section{Per-seed distributions}
\label{app:seed_distributions}

\Cref{tab:fixed17_alignment} reports arithmetic means across retained systems after per-system seed-IQM summaries. The corresponding per-(system, seed) distributions are shown below. \Cref{fig:fixed17_seed_strips} shows the paper-facing strip plot for the matched boundary-emphasized rows. \Cref{fig:appfig_fixed17_perseed} plots per-(system, seed) histograms of forecasting MSE at every paper-facing horizon. \Cref{fig:appfig_fixed17_entropy} plots per-(system, seed) histograms of \(H(B\mid S_{\rm abs})\), \(H(S_{\rm abs}\mid B)\), and exact-support uniqueness on the deep-state slice. \Cref{fig:appfig_dysts_perseed} plots the same per-seed distributions for the Dysts \(dt{\times}30\) benchmark. The histograms show retained-system and seed-level dispersion around the paper-facing means.

\begin{figure}[tbp]
\centering
\includegraphics[width=\linewidth]{neurips_paper_2026/fig_fixed17_seed_strips_boundary_only.pdf}
\caption{Per-(system, seed) forecasting MSE at \(H100\), \(H500\), and \(H1000\) on the multibasin benchmark. Black bars mark the mean across per-system seed-IQM values. The Dense MLP has a much heavier upper tail at every horizon.}
\label{fig:fixed17_seed_strips}
\end{figure}

\begin{figure}[tbp]
\centering
\includegraphics[width=\linewidth]{neurips_paper_2026/appfig_fixed17_perseed_histograms.pdf}
\caption{Per-(system, seed) forecasting MSE distributions on the multibasin benchmark. Each row is a horizon; histograms are over the \(\log_{10}\) raw MSE.}
\label{fig:appfig_fixed17_perseed}
\end{figure}

\begin{figure}[tbp]
\centering
\includegraphics[width=\linewidth]{neurips_paper_2026/appfig_fixed17_entropy_histograms.pdf}
\caption{Per-(system, seed) support--basin diagnostic distributions on the deep-state slice. Histograms show \(H(B\mid S_{\rm abs})\), \(H(S_{\rm abs}\mid B)\), and exact-support uniqueness for the rows summarized in \cref{tab:fixed17_alignment}.}
\label{fig:appfig_fixed17_entropy}
\end{figure}

\begin{figure}[tbp]
\centering
\includegraphics[width=\linewidth]{neurips_paper_2026/appfig_dysts_dt30_perseed_histograms.pdf}
\caption{Per-(system, seed) forecasting MSE distributions on the Dysts \(dt{\times}30\) benchmark.}
\label{fig:appfig_dysts_perseed}
\end{figure}

\section{Support-routing robustness}
\label{app:routing_robustness}

\Cref{tab:self_routing_deep} reports route-form diagnostics that are not the main table. The deep slice removes boundary-adjacent states, so it is useful as a robustness check rather than the primary deployment setting. The exact \(S_{\rm top8}\) columns should be read as ablations of the support-family router rather than as the headline routing object.

For the gated exact-support diagnostic, let \(m_c\in\{0,1\}^{d_z}\) be the eight-coordinate mask of the exact route \(c=S_{\rm top8}(z_t)\), and let \(\bar z_c\) be the corresponding fitting-set mean. The gated rule masks the centered latent state before applying the learned global transition:
\begin{equation}
  \hat z_{t+1}=\bar z_c+K\bigl(m_c\odot (z_t-\bar z_c)\bigr).
  \label{eq:gated_k}
\end{equation}
This diagnostic asks whether the active coordinates select useful columns or subspaces of the learned transition. The family-local route in the main text asks a different question: whether the same support family can key a post-hoc local linear map.

\begin{table}[tbp]
\centering
\caption{Deep-state support-routed forecasting route-form diagnostics. \(S_{\rm top8}\) columns use exact eight-index supports; \(F_{\rm top8}\) columns use merged families of those exact supports. Entries are arithmetic means across systems after within-system seed-IQM summarization, with diagnostic within-system Wilcoxon/Holm counts in brackets. LISTA-SB is the matched \(d_z{=}256\) soft-block LISTA ablation; lower is better.}
\label{tab:self_routing_deep}
\begin{minipage}[t]{0.49\textwidth}
\centering
\textbf{\(H100\)}
\resizebox{\linewidth}{!}{
\begin{tabular}{@{}l rrr@{}}
\toprule
Model & \shortstack{Gated \(S_{\rm top8}\)\\\cref{eq:gated_k}} & \shortstack{\(S_{\rm top8}\)-local\\\cref{eq:local_fit}} & \shortstack{\(F_{\rm top8}\)-local\\\cref{eq:local_fit}} \\
\midrule
LISTA          & \(0.783\,[4/15]\)           & \(0.771\,[5/15]\)           & \(0.165\,[9/15]\) \\
LISTA-SB       & \(0.861\,[1/15]\)           & \(0.861\,[1/15]\)           & \(0.00372\,[13/15]\) \\
LISTA-BD       & \(0.915\,[1/15]\)           & \(0.919\,[1/15]\)           & \(\mathbf{0.00207}\,[12/15]\) \\
Sparse MLP, BD & \(\mathbf{0.797}\,[0/15]\) & \(0.770\,[0/15]\)           & \(0.00941\,[6/15]\) \\
Sparse MLP     & \(0.813\,[0/15]\)           & \(\mathbf{0.736}\,[0/15]\) & \(0.0104\,[5/15]\) \\
Dense MLP      & \(0.991\,[0/15]\)           & \(0.990\,[0/15]\)           & \(102\,[0/15]\) \\
\bottomrule
\end{tabular}
}
\end{minipage}\hfill
\begin{minipage}[t]{0.49\textwidth}
\centering
\textbf{\(H1000\)}
\resizebox{\linewidth}{!}{
\begin{tabular}{@{}l rrr@{}}
\toprule
Model & \shortstack{Gated \(S_{\rm top8}\)\\\cref{eq:gated_k}} & \shortstack{\(S_{\rm top8}\)-local\\\cref{eq:local_fit}} & \shortstack{\(F_{\rm top8}\)-local\\\cref{eq:local_fit}} \\
\midrule
LISTA          & \(0.681\,[8/15]\)           & \(0.676\,[7/15]\)           & \(\mathbf{8.03{\times}10^{-14}}\,[11/15]\) \\
LISTA-SB       & \(0.817\,[7/15]\)           & \(0.809\,[6/15]\)           & \(6.30{\times}10^{-11}\,[13/15]\) \\
LISTA-BD       & \(0.690\,[6/15]\)           & \(0.709\,[7/15]\)           & \(1.19{\times}10^{-7}\,[13/15]\) \\
Sparse MLP, BD & \(\mathbf{1.48{\times}10^{-5}}\,[1/15]\) & \(\mathbf{2.99{\times}10^{-4}}\,[1/15]\) & \(1.49{\times}10^{-6}\,[14/15]\) \\
Sparse MLP     & \(0.804\,[2/15]\)           & \(0.738\,[1/15]\)           & \(0.126\,[14/15]\) \\
Dense MLP      & \(0.849\,[0/15]\)           & \(0.965\,[0/15]\)           & \(78.6\,[0/15]\) \\
\bottomrule
\end{tabular}
}
\end{minipage}
\end{table}

\section{Dysts forecasting diagnostics}
\label{app:dysts_full}

The per-(system, seed) Dysts MSE values are released in the collected \texttt{forecasting\_rows.csv} artifacts. \Cref{tab:dysts_long} reports arithmetic means across systems after first summarizing seeds within each system by IQM. The \(dt{\times}30\) packet is fully finite in median coverage at every reported horizon, but the per-seed distributions remain heavy-tailed on several systems, so the main table and \cref{fig:dysts_long_horizon,fig:appfig_dysts_perseed} should be read together: the table gives actual-MSE point estimates and paired system-level significance markers, while the figures show horizon-scale uncertainty and per-seed dispersion. The soft-block LISTA diagnostic improves on Dense MLP in aggregate but does not match the primary sparse forecasting rows and does not survive the main Dysts significance correction. The older tiny-step \(H{\leq}60000\) packet is useful as a stress-test diagnostic because it exposed divergent block-diagonal LISTA rollouts, but it is not pooled with the coarser-step table.

\begin{table}[tbp]
\centering
\caption{Scale-normalized Dysts forecasting summary. Entries are arithmetic means across systems of the candidate per-system seed-IQM MSE divided by the Dense MLP per-system seed-IQM MSE; lower than \(1\) favors the candidate. This table gives the average-case ratio view corresponding to the actual-MSE table in \cref{tab:dysts_long}. LISTA-SB is retained as a diagnostic row.}
\label{tab:dysts_ratio_to_dense}
\resizebox{\textwidth}{!}{%
\begin{tabular}{@{}l r r r r r r r r@{}}
\toprule
Model & H100 & H500 & H1000 & H1500 & H2000 & H3000 & H4000 & H5000 \\
\midrule
LISTA & \(0.509\) & \(0.464\) & \(0.459\) & \(0.426\) & \(0.440\) & \(0.489\) & \(0.565\) & \(0.595\) \\
LISTA-BD & \(0.530\) & \(0.587\) & \(0.502\) & \(0.445\) & \(0.443\) & \(0.498\) & \(0.574\) & \(0.612\) \\
LISTA-SB & \(0.972\) & \(1.026\) & \(1.019\) & \(1.088\) & \(1.157\) & \(1.264\) & \(1.086\) & \(1.006\) \\
Sparse MLP & \(0.462\) & \(0.499\) & \(0.501\) & \(0.563\) & \(0.566\) & \(0.610\) & \(0.685\) & \(0.739\) \\
Sparse MLP-BD & \(0.477\) & \(0.499\) & \(0.497\) & \(0.555\) & \(0.559\) & \(0.617\) & \(0.703\) & \(0.757\) \\
Dense MLP & \(1.000\) & \(1.000\) & \(1.000\) & \(1.000\) & \(1.000\) & \(1.000\) & \(1.000\) & \(1.000\) \\
\bottomrule
\end{tabular}
}
\end{table}

\section{Additional background}
\label{app:background}

\subsection{Koopman operators and finite-dimensional approximations}
\label{app:koopman_background}

Consider the stored-step map used in \cref{sec:problem},
\begin{equation}
  x_{k+1}=F_{\Delta t}(x_k), \qquad x_k\in\mathcal X\subseteq\R^{d_x}.
\end{equation}
Koopman operator theory studies the evolution of observables rather than the evolution of the state directly. An observable is a measurement function \(g:\mathcal X\to\R\), and the Koopman operator \(\mathcal K\) acts by composition:
\begin{equation}
  (\mathcal K g)(x)=g(F_{\Delta t}(x)).
\end{equation}
The map \(F_{\Delta t}\) may be nonlinear, but \(\mathcal K\) is linear on the function space of observables. If a finite vector of observables
\begin{equation}
  \Phi(x)=\big[g_1(x),\ldots,g_{d_z}(x)\big]^\top
\end{equation}
spans a Koopman-invariant subspace, then there is a matrix \(K\in\R^{d_z\times d_z}\) such that
\begin{equation}
  \Phi(x_{k+1})=\Phi(F_{\Delta t}(x_k))=\mathcal K\Phi(x_k)=K\Phi(x_k).
\end{equation}
In that case the nonlinear dynamics are exactly linear in the observable coordinates \(z_k=\Phi(x_k)\). Learning Koopman models can therefore be viewed as searching for observables whose span is invariant enough to support useful finite-dimensional linear prediction.

In practice, the infinite-dimensional operator is approximated from data. Given snapshot pairs \(\{(x_k,x_{k+1})\}\) and a chosen feature map \(\Phi\), dynamic mode decomposition and extended dynamic mode decomposition (eDMD) estimate a finite transition by least squares \citep{schmid_dmd_2010,rowley_2009_spectral,tu_dynamic_2014,williams_datadriven_2015,williams_kernel-based_2016}:
\begin{equation}
  K
  =
  \argmin_{\widetilde K\in\R^{d_z\times d_z}}
  \sum_k \norm{\Phi(x_{k+1})-\widetilde K\Phi(x_k)}_2^2.
\end{equation}
DMD is the special case \(\Phi(x)=x\), so it fits a linear state-space surrogate without a learned nonlinear lift. eDMD extends this by choosing a richer observable dictionary. Koopman autoencoders replace the fixed observable dictionary with a learned encoder and decoder, which is the setting used in this paper.

\subsection{Why multibasin systems may benefit from local structure}
\label{app:multibasin_background}

Near appropriate attracting sets, classical linearization and Koopman eigenfunction constructions can produce useful local or basin-restricted coordinates. Global results are more restrictive: multi-attractor systems should not in general be assumed to admit one exact finite-dimensional linearization with a continuous, state-reconstructing encoder--decoder pair \citep{lan_linearization_2013,brunton_koopman_2016,pan_lifting_2024,liu_properties_2025}. These results motivate local structure, but they do not prove that every learned global KAE must fail or that sparse supports recover the theoretical local coordinates.

For the present paper, a learned finite-dimensional KAE is treated as an approximation. Sparse supports provide one possible discrete and inspectable key: a support may index a basin part or local predictive chart while coefficient values carry within-chart information. Whether that interpretation holds is an empirical question tested by the transductive alignment and staged forecasting experiments.

\subsection{Koopman autoencoders and periodic re-encoding}
\label{app:kae_reencoding_background}

A standard Koopman autoencoder learns an encoder \(\Enc:\R^{d_x}\to\R^{d_z}\), decoder \(\Dec:\R^{d_z}\to\R^{d_x}\), and latent transition \(K\) so that
\begin{equation}
  z_k=\Enc(x_k),\qquad z_{k+1}\approx Kz_k,\qquad x_k\approx\Dec(z_k).
\end{equation}
A schematic one-step objective has the form
\begin{equation}
  \min_{\Enc,\Dec,K}
  \sum_k
  \norm{x_k-\Dec(\Enc(x_k))}_2^2
  +
  \lambda_{\rm lin}
  \norm{\Enc(x_{k+1})-K\Enc(x_k)}_2^2 .
  \label{eq:KAE-standard-objective}
\end{equation}
This objective is useful as a reference point, but it is not sufficient for the present study. It does not directly train decoded multi-step rollouts, does not impose a sparse support object, and does not test whether the support itself selects a useful local predictive rule. The training objective in \cref{sec:method_training} therefore keeps the autoencoding and latent-linearity terms but adds decoded rollout prediction over windows, sparsity on the rolled latent states, and optional transition-structure penalties.

Periodic re-encoding interrupts a latent rollout after a fixed number of stored observation steps. For period \(m\),
\begin{equation}
  \tilde z_{k+m}=K^m z_k,\qquad
  \tilde x_{k+m}=\Dec(\tilde z_{k+m}),\qquad
  z_{k+m}=\Enc(\tilde x_{k+m}).
\end{equation}
This is the same mechanism defined in \cref{eq:reencode}. It uses only the model's own predicted state and is therefore not a teacher-forced correction. The point is to refresh the latent embedding and, in sparse models, to refresh the support as the predicted state moves through regions where a different local linear description may be more appropriate \citep{fathi2024course}.

\subsection{Sparse coding and LISTA}
\label{app:sparse_lista_background}

Classical sparse coding seeks a code \(z\) that reconstructs a signal from a small set of active dictionary atoms:
\begin{equation}
  z^\star(x)=
  \argmin_{z\in\R^{d_z}}\;\frac{1}{2}\norm{x-Dz}_2^2+\lambda_{\rm sc}\norm{z}_1,
  \qquad \lambda_{\rm sc}>0,
\end{equation}
where \(D\) is a dictionary and \(\lambda_{\rm sc}\) controls sparsity \citep{olshausen_emergence_1996,chen_atomic_1998,donoho_optimally_2003}. For a fixed dictionary, and under the usual uniqueness and non-boundary conditions for the Lasso solution, the active set and signs are locally constant over regions of the input space, while the coefficient values vary continuously within a fixed active set. The decoder reconstruction then lies in the span of the selected dictionary columns, giving a union-of-subspaces geometry. Each exact support can index a local chart, and the greedy Jaccard procedure in \cref{sec:method_supports} groups nearby Boolean support vectors into support families when threshold noise or boundary effects fragment the exact masks.

LISTA replaces an iterative sparse-coding solver with a learned finite-depth shrinkage network \citep{gregor_learning_2010}. In the notation of \cref{eq:lista}, the encoder forms a pre-code from the input and repeatedly applies learned affine refinement followed by soft thresholding. In this paper, the encoder is trained jointly with the Koopman rollout objective rather than solving the sparse-coding problem at evaluation time. Thresholding turns the latent into continuous coefficients plus an explicit support. The experiments ask whether absolute-threshold support families align transductively with evaluation-only native or nearest-center-proxy labels and whether a separately fitted training-distribution family partition can route useful local affine predictors.


\bibliographystyle{plainnat}
\bibliography{neurips_sparse_koopman_multibasin}

\end{document}
